% Standalone section on site-level primary-source identification.
% Compile with: pdflatex main.tex
% Drop into the manuscript with % Standalone section on site-level primary-source identification.
% Compile with: pdflatex main.tex
% Drop into the manuscript with % Standalone section on site-level primary-source identification.
% Compile with: pdflatex main.tex
% Drop into the manuscript with % Standalone section on site-level primary-source identification.
% Compile with: pdflatex main.tex
% Drop into the manuscript with \input{sec_primary_driver}
% after defining nothing: the section is self-contained.

\section{Primary-source identification across heterogeneous sites}
\label{sec:primary-source}

The manuscript separates covariate shift, a change in $P(X)$, from concept
drift, a change in $P(Y\mid X)$ between a reference batch
$D_{\mathrm{exist}}$ and a new batch $D_{\mathrm{new}}$.
Feature-level procedures in the paper, including the pseudo-outcome
risk and conditional permutation importance, score coordinates inside one
batch.
The federated aggregation in Appendix~11 then asks whether those
variable-importance vectors can be merged without moving raw rows.
That merge treats every site as a noisy look at one common shift.
A prior operational question is which sites are the primary sources of the
shift, and which sites are either stable or actively destroying the
consensus.
This section constructs a ten-site benchmark in which the answer is known,
and evaluates a counterfactual concept score against that answer.

The design isolates concept drift.
Every site is a random slice of one real covariate table, so the marginal
distribution $P(X)$ is shared.
The American Community Survey extracts below are cross-sections, not
panels, so the batch index is a seeded random order rather than a clock.
Stationarity of $P(X)$ is then a construction, and a covariate arm that
fires is evidence of a bug in the generator rather than a feature of the
task.
Labels follow the real outcome through a burn-in and are then rewritten,
on a known subset of sites, by a known rule.

\subsection{Data and the ten-site generator}
\label{sec:dgp}

Two WhyShift tasks supply the base table $(X,y^{\mathrm{base}})$, both
taken from the 2018 ACS one-year person file for California and both using
the official \texttt{folktables} feature lists.
ACS Income has $p=10$ covariates and the indicator that personal income
exceeds $\$50{,}000$.
ACS Public Coverage has $p=19$ covariates and the indicator of public
health-insurance coverage.
In each Monte Carlo replicate the analysis sample is winsorized at the
first and ninety-ninth percentiles, standardized, and cut into $K=10$
non-overlapping sites of $n_k=1600$ rows.
Each site is then cut into $T=16$ batches of $100$ rows.
The burn-in is the first four batches.
The concept change, when it is switched on, begins at batch
$t^\star=8$.

The ten sites have four roles.
Four clean sites keep $y=y^{\mathrm{base}}$ at every batch.
Three linear sites rewrite labels inside an upper tail of a designated
coordinate $X_1$.
Two nonlinear sites rewrite every activated row.
One noise site replaces activated labels with independent
$\mathrm{Bernoulli}(1/2)$ draws.
The noise site is a consensus breaker.
It is not a primary source: there is no new concept for a local model to
learn.
Table~\ref{tab:roles} records which coordinates carry the rules.
Income uses hours worked, educational attainment, and age.
Public Coverage uses age, educational attainment, and personal income.
These are the continuous or ordinal columns in each official feature list.
The remaining columns stay in the table as nuisance covariates.

\begin{table}[t]
\centering
\caption{Rule coordinates. $Z_j$ denotes the analysis-sample standardize
of the named column. Linear sites rewrite only the upper tail of $Z_1$.
Nonlinear sites rewrite every activated row.}
\label{tab:roles}
\begin{tabular}{llccc}
\toprule
Task & Role of $(X_1,X_2,X_3)$ & $X_1$ & $X_2$ & $X_3$ \\
\midrule
ACS Income & hours, schooling, age & WKHP & SCHL & AGEP \\
ACS Public Coverage & age, schooling, income & AGEP & SCHL & PINCP \\
\bottomrule
\end{tabular}
\end{table}

Write $Z$ for the standardized covariates and let
$M_k(q)=\{Z_1>Q_q(Z_1)\}$ be the upper tail at quantile $q$, computed on
the pooled analysis sample.
The main design uses $q=0.5$.
On a linear site the candidate label is the halfspace
\begin{equation}
\label{eq:linear-rule}
y^{\mathrm{rule}}
=
\mathbf{1}\!\left\{
\sigma_k(\beta_{k1} Z_1+\beta_{k2} Z_2)>c_k
\right\}
\quad\text{on }M_k(q),
\end{equation}
and $y^{\mathrm{base}}$ off that tail.
The three coefficient vectors are
$(1,0.35)$, $(0.25,1)$, and $(-0.85,0.55)$.
The threshold $c_k$ is the quantile of the linear score that matches the
positive rate of $y^{\mathrm{base}}$ on $M_k(q)$, and the sign
$\sigma_k\in\{\pm 1\}$ is the orientation whose matched labels agree less
with $y^{\mathrm{base}}$.
The sign choice is part of the generator.
It keeps the planted concept from collapsing back onto the real outcome
when hours or schooling already predict income.
It is not estimated by the detector.

On a nonlinear site the same rate match is applied to every row, not only
the tail:
\begin{equation}
\label{eq:nonlinear-rule}
y^{\mathrm{rule}}
=
\mathbf{1}\{Z_1 Z_2>c\}
\qquad\text{or}\qquad
y^{\mathrm{rule}}
=
\mathbf{1}\{\sin(Z_1)+Z_3^2>c\}.
\end{equation}
The first site uses the product and the second uses the sine-square rule.
The planted feature support is $\{X_1,X_2\}$ for every linear site and for
the product site, and $\{X_1,X_3\}$ for the sine-square site.

Activation is a batch weight $w_t\in[0,1]$.
Under abrupt activation, $w_t=\mathbf{1}\{t\ge t^\star\}$.
Under gradual activation, $w_t=0$ before $t^\star$ and then rises linearly
to one over four batches.
Each post-burn-in row is rewritten with probability $\alpha w_t$, and
otherwise keeps $y^{\mathrm{base}}$.
The main grid uses $\alpha=1$.
The noise site uses the same weights and then draws a fresh coin flip.
Clean sites ignore the weight.

Two consequences are worth stating before any score is computed.
First, because the structured rules match the local positive rate, a
marginal label discrepancy is not a substitute for concept drift.
The largest average gap between the post-change positive rate and the
baseline, among structured sites, is $2.4$ percentage points
(Income, nonlinear).
Second, hours worked in the ACS piles up at $40$.
The empirical quantiles of WKHP at $0.5$ and at $0.7$ are both $40$, so
$\{Z_1>Q_{0.5}\}$ and $\{Z_1>Q_{0.7}\}$ are the same set, about $23\%$ of
Income rows.
The median tail in the main Income design is therefore thinner than a
continuous coordinate would suggest.
Public Coverage uses age as $X_1$, and that tail is about half the sample.
Table~\ref{tab:dgp-check} records the realized flip rates.
Structured sites move roughly half the labels inside the region they
rewrite, and they leave the marginal positive rate near the baseline.
The noise site flips about half of all post-burn-in labels and shifts the
positive rate to one half.

\begin{table}[t]
\centering
\caption{Realized label changes after $t^\star$, averaged over four seeds
at $\alpha=1$, abrupt activation, and $q=0.5$.
Flip rate is the fraction of post-burn-in rows whose label differs from
$y^{\mathrm{base}}$.
Agreement is the fraction of the rewritten region on which the rate-matched
rule reproduces $y^{\mathrm{base}}$.
Clean sites are not rewritten, so agreement is one.
The noise site has no rule, so agreement is left blank.}
\label{tab:dgp-check}
\begin{tabular}{llcccc}
\toprule
Task & Site type & Flip rate & Rule agreement & Post positive rate & Base positive rate \\
\midrule
Income & Clean & 0.000 & 1.000 & 0.406 & 0.406 \\
Income & Linear & 0.131 & 0.427 & 0.400 & 0.414 \\
Income & Nonlinear & 0.555 & 0.436 & 0.396 & 0.420 \\
Income & Noise & 0.498 & --- & 0.491 & 0.410 \\
Public Coverage & Clean & 0.000 & 1.000 & 0.369 & 0.367 \\
Public Coverage & Linear & 0.259 & 0.476 & 0.378 & 0.372 \\
Public Coverage & Nonlinear & 0.492 & 0.506 & 0.371 & 0.374 \\
Public Coverage & Noise & 0.496 & --- & 0.495 & 0.362 \\
\bottomrule
\end{tabular}
\end{table}

\subsection{The counterfactual concept score}
\label{sec:delta-cf}

Let $\mathcal{D}_0$ be the pooled burn-in, eight thousand rows in the main
design, still generated from $y^{\mathrm{base}}$ at every site.
Let $\mathcal{D}_k$ be the most recent six batches at site $k$, six hundred
rows, all of them inside the drift window when activation is abrupt.
A random forest with $25$ trees, depth $8$, and minimum leaf size $5$ is
the outcome model throughout, in line with the recommendation that the
forest be used for distance estimation rather than for a causal reading of
the meta-learner.
Local predictions on $\mathcal{D}_k$ are three-fold cross-fit probabilities
$\mu_k^{\mathrm{oof}}$.
The reference probabilities $\mu_0^{\mathrm{sm}}$ are an average of two
forests, each trained on a random subset of $\mathcal{D}_0$ whose size
equals the local training set, two thirds of the current window.
Matching the training size stops the larger reference sample from looking
like a concept change.
A third forest trained on all of $\mathcal{D}_0$ is retained only as the
pooled excess defined below.
It is not the headline score.

On a region $\mathcal{R}$ that has both classes in the current window and
in the reference, at least fifty current rows, and at least eighty
reference rows, write $B(y,p)$ for the Brier score and
\begin{equation}
\label{eq:delta-region}
\Delta_k(\mathcal{R})
=
\Big(B(\mu_0^{\mathrm{sm}})-B(\mu_k^{\mathrm{oof}})\Big)_+
\cdot
\Big(2\big(\mathrm{AUC}(\mu_k^{\mathrm{oof}})-{\textstyle\frac12}\big)\Big)_+.
\end{equation}
The first factor is the excess risk of the reference concept over a local
concept of the same capacity.
The second factor is zero when the local model does not beat a coin, and
it equals one when the local AUC equals one.
The headline score is the maximum over a dictionary that does not use
labels,
\begin{equation}
\label{eq:delta-cf}
\Delta_k^{\mathrm{cf}}
=
\max_{\mathcal{R}\in\{\mathrm{all},\,\mathrm{high},\,\mathrm{low}\}}
\Delta_k(\mathcal{R}),
\end{equation}
where ``high'' and ``low'' are the two sides of the burn-in median of
$X_1$.
The median is computed on $\mathcal{D}_0$ only.
The whole-window score $\Delta_k(\mathrm{all})$ is reported beside the
maximum so that the cost of ignoring the tail is visible.

The same two probability vectors define the concept contrast
\begin{equation}
\label{eq:tau}
\tau_k^2
=
\frac{1}{|\mathcal{D}_k\cap\mathcal{R}_k^\star|}
\sum_{i\in\mathcal{D}_k\cap\mathcal{R}_k^\star}
\big(\mu_k^{\mathrm{oof}}(X_i)-\mu_0^{\mathrm{sm}}(X_i)\big)^2,
\end{equation}
with $\mathcal{R}_k^\star$ the region attaining the maximum in
\eqref{eq:delta-cf}.
This is the concept contrast that the pseudo-outcome decomposition isolates
from the covariate arm once $P(X)$ is shared across $D_{\mathrm{exist}}$
and $D_{\mathrm{new}}$.
The benchmark does not re-estimate a propensity: under the common
covariate table that weight is not where the signal sits.
Gating by the same skill factor as in \eqref{eq:delta-region} produces
$\tau_k^{2,\mathrm{g}}$.
Raw $\tau_k^2$ is retained because a flat local predictor on a noise site
can still sit far from a sharp reference predictor.
The gate is the part of $\Delta_k^{\mathrm{cf}}$ that refuses to call that
gap a primary source.

For a score $s$ and two disjoint collections of sites $\mathcal{A}$ and
$\mathcal{B}$, the pairwise concordance is
\begin{equation}
\label{eq:concordance}
C(s;\mathcal{A},\mathcal{B})
=
\frac{1}{|\mathcal{A}|\,|\mathcal{B}|}
\sum_{a\in\mathcal{A}}\sum_{b\in\mathcal{B}}
\mathbf{1}\{s_a>s_b\}.
\end{equation}
Ties count as failures.
The primary set $\mathcal{P}$ is the five structured sites.
The target for the headline score is
$C(\Delta^{\mathrm{cf}};\mathcal{P},\mathcal{C})>0.8$,
where $\mathcal{C}$ is the set of four clean sites.
The same functional is reported for $\mathcal{P}$ against the noise site,
and separately for the linear and nonlinear subsets of $\mathcal{P}$.
Every entry below is a mean over four seeds, with the worst seed in
parentheses, unless a row is explicitly a per-type mean.

Negative controls use the same windows and never see the label rule.
The domain AUC is the cross-fit AUC of a classifier trained to separate
burn-in rows from recent rows using $X$ alone, once with the forest and
once with logistic regression.
The covariate discrepancy is an RBF MMD on subsamples of size $250$, with
bandwidth equal to the median of the squared Euclidean distances.
A joint classifier two-sample test appends $y$ to $X$ before the same
forest AUC is computed.
Marginal label total variation is the absolute difference of positive rates
between the recent window and the burn-in.

\begin{table}[t]
\centering
\caption{Mean score by site type at $\alpha=1$, abrupt activation, and
$q=0.5$. Four seeds. $\Delta^{\mathrm{cf}}$ is \eqref{eq:delta-cf}.
$\Delta(\mathrm{all})$ is the whole-window score.
$\Delta(\mathrm{high})$ is the upper tail of $X_1$.
Domain AUC is the logistic classifier on $X$ only.}
\label{tab:means}
\small
\begin{tabular}{llrrrrrr}
\toprule
Task & Type & $\Delta^{\mathrm{cf}}$ & $\Delta(\mathrm{all})$ & $\Delta(\mathrm{high})$ & $\tau^2$ & Skill & Domain AUC \\
\midrule
Income & Clean & 0.003 & 0.000 & 0.002 & 0.015 & 0.696 & 0.498 \\
Income & Linear & 0.253 & 0.024 & 0.253 & 0.181 & 0.937 & 0.482 \\
Income & Nonlinear & 0.336 & 0.298 & 0.316 & 0.313 & 0.897 & 0.494 \\
Income & Noise & 0.001 & 0.000 & 0.000 & 0.096 & 0.008 & 0.494 \\
Public Coverage & Clean & 0.001 & 0.000 & 0.001 & 0.014 & 0.364 & 0.503 \\
Public Coverage & Linear & 0.200 & 0.052 & 0.200 & 0.132 & 0.943 & 0.493 \\
Public Coverage & Nonlinear & 0.236 & 0.182 & 0.179 & 0.159 & 0.970 & 0.478 \\
Public Coverage & Noise & 0.003 & 0.002 & 0.003 & 0.046 & 0.089 & 0.528 \\
\bottomrule
\end{tabular}
\end{table}

Table~\ref{tab:means} is the raw material of the concordance claim.
On both tasks, and for every linear site, the maximum is exactly the
upper-tail score: $\Delta^{\mathrm{cf}}=\Delta(\mathrm{high})$.
On Income that tail score is $0.253$ against a whole-window score of
$0.024$.
The whole window is mostly unrewritten rows, and a single Brier average
dilutes a real halfspace by about an order of magnitude.
The noise skill is $0.008$ on Income and $0.089$ on Public Coverage, so
the gate shuts even though the ungated contrast $\tau^2$ on Income noise
($0.096$) exceeds the clean contrast ($0.015$).
Logistic domain AUCs sit on $1/2$ for every type.
The largest is $0.528$, on the Public Coverage noise site.
The forest domain AUC, not shown column by column, has type means between
$0.490$ and $0.514$, and equals $0.514$ on that same noise site.
Neither is evidence of a covariate shift.

\subsection{Concordance}
\label{sec:concordance}

\begin{table}[t]
\centering
\caption{Pairwise concordance \eqref{eq:concordance} at $\alpha=1$, abrupt
activation, and $q=0.5$.
Each cell is the mean over four seeds, with the worst seed in parentheses.
$\mathcal{P}$ is the five structured sites and $\mathcal{C}$ is the four
clean sites.}
\label{tab:concordance}
\small
\begin{tabular}{llcccc}
\toprule
Task & Score & $\mathcal{P}>\mathcal{C}$ & $\mathcal{P}>\mathrm{noise}$ & Linear $>\mathcal{C}$ & Nonlinear $>\mathcal{C}$ \\
\midrule
Income & $\Delta^{\mathrm{cf}}$ & 1.000 {\scriptsize(1.000)} & 1.000 {\scriptsize(1.000)} & 1.000 {\scriptsize(1.000)} & 1.000 {\scriptsize(1.000)} \\
Income & $\Delta(\mathrm{all})$ & 1.000 {\scriptsize(1.000)} & 1.000 {\scriptsize(1.000)} & 1.000 {\scriptsize(1.000)} & 1.000 {\scriptsize(1.000)} \\
Income & $\tau^{2,\mathrm{g}}$ & 1.000 {\scriptsize(1.000)} & 1.000 {\scriptsize(1.000)} & 1.000 {\scriptsize(1.000)} & 1.000 {\scriptsize(1.000)} \\
Income & Pareto rank & 1.000 {\scriptsize(1.000)} & 1.000 {\scriptsize(1.000)} & 1.000 {\scriptsize(1.000)} & 1.000 {\scriptsize(1.000)} \\
Income & $\tau^2$ & 1.000 {\scriptsize(1.000)} & 1.000 {\scriptsize(1.000)} & 1.000 {\scriptsize(1.000)} & 1.000 {\scriptsize(1.000)} \\
Income & Pooled excess & 1.000 {\scriptsize(1.000)} & 0.400 {\scriptsize(0.400)} & 1.000 {\scriptsize(1.000)} & 1.000 {\scriptsize(1.000)} \\
Income & Label TV & 0.500 {\scriptsize(0.100)} & 0.050 {\scriptsize(0.000)} & 0.500 {\scriptsize(0.083)} & 0.500 {\scriptsize(0.125)} \\
Income & Joint C2ST & 0.850 {\scriptsize(0.800)} & 0.350 {\scriptsize(0.200)} & 0.750 {\scriptsize(0.667)} & 1.000 {\scriptsize(1.000)} \\
Income & RF domain AUC & 0.487 {\scriptsize(0.400)} & 0.500 {\scriptsize(0.000)} & 0.521 {\scriptsize(0.333)} & 0.438 {\scriptsize(0.250)} \\
Income & RBF MMD & 0.412 {\scriptsize(0.300)} & 0.500 {\scriptsize(0.000)} & 0.375 {\scriptsize(0.167)} & 0.469 {\scriptsize(0.125)} \\
\midrule
Public Coverage & $\Delta^{\mathrm{cf}}$ & 1.000 {\scriptsize(1.000)} & 1.000 {\scriptsize(1.000)} & 1.000 {\scriptsize(1.000)} & 1.000 {\scriptsize(1.000)} \\
Public Coverage & $\Delta(\mathrm{all})$ & 1.000 {\scriptsize(1.000)} & 1.000 {\scriptsize(1.000)} & 1.000 {\scriptsize(1.000)} & 1.000 {\scriptsize(1.000)} \\
Public Coverage & $\tau^{2,\mathrm{g}}$ & 1.000 {\scriptsize(1.000)} & 1.000 {\scriptsize(1.000)} & 1.000 {\scriptsize(1.000)} & 1.000 {\scriptsize(1.000)} \\
Public Coverage & Pareto rank & 1.000 {\scriptsize(1.000)} & 1.000 {\scriptsize(1.000)} & 1.000 {\scriptsize(1.000)} & 1.000 {\scriptsize(1.000)} \\
Public Coverage & $\tau^2$ & 1.000 {\scriptsize(1.000)} & 1.000 {\scriptsize(1.000)} & 1.000 {\scriptsize(1.000)} & 1.000 {\scriptsize(1.000)} \\
Public Coverage & Pooled excess & 1.000 {\scriptsize(1.000)} & 1.000 {\scriptsize(1.000)} & 1.000 {\scriptsize(1.000)} & 1.000 {\scriptsize(1.000)} \\
Public Coverage & Label TV & 0.613 {\scriptsize(0.350)} & 0.000 {\scriptsize(0.000)} & 0.604 {\scriptsize(0.417)} & 0.625 {\scriptsize(0.250)} \\
Public Coverage & Joint C2ST & 0.950 {\scriptsize(0.850)} & 0.450 {\scriptsize(0.200)} & 0.917 {\scriptsize(0.750)} & 1.000 {\scriptsize(1.000)} \\
Public Coverage & RF domain AUC & 0.438 {\scriptsize(0.350)} & 0.200 {\scriptsize(0.000)} & 0.396 {\scriptsize(0.167)} & 0.500 {\scriptsize(0.250)} \\
Public Coverage & RBF MMD & 0.600 {\scriptsize(0.500)} & 0.200 {\scriptsize(0.000)} & 0.562 {\scriptsize(0.417)} & 0.656 {\scriptsize(0.375)} \\
\bottomrule
\end{tabular}
\end{table}

Table~\ref{tab:concordance} is the primary result.
On both tasks and in every seed,
$C(\Delta^{\mathrm{cf}};\mathcal{P},\mathcal{C})=1$,
and the same equality holds for linear sites alone, for nonlinear sites
alone, and for $\mathcal{P}$ against the noise site.
The whole-window score and the skill-gated contrast $\tau^{2,\mathrm{g}}$
match this ranking at full intensity.
Figure~\ref{fig:concordance} plots the clean-site comparison.
The covariate rows do not.
Random-forest domain AUC and RBF MMD have concordances between $0.38$ and
$0.66$, with worst seeds as low as $0.17$.
Label total variation is consistent with a matched positive rate: its
concordance against clean sites is $0.50$ on Income and $0.61$ on Public
Coverage, and its concordance against the noise site is $0.05$ and $0.00$.
The noise site is the unique marginal shift, and a marginal functional
ranks it first.

The joint classifier two-sample test sees any change in $(X,Y)$.
It separates $\mathcal{P}$ from $\mathcal{C}$ at $0.85$ on Income and $0.95$
on Public Coverage, which clears $0.8$ on average, but it ranks the noise
site as the stronger departure: concordance of $\mathcal{P}$ against noise
is $0.35$ and $0.45$.
On Income the linear-only comparison for the joint test falls to $0.75$.
A test of the joint distribution is not a test for a learnable new concept.

Pooled excess makes the same mistake in a sharper way, and only on the
task where the reference concept is easy.
It is the Brier gap between the forest trained on the entire burn-in and
the local cross-fit model, before size matching and before the skill gate.
On Income its concordance against the noise site is $0.400$ in every seed.
That fraction is exact: both nonlinear sites beat the noise site, and none
of the three linear sites do.
The pooled Income forest is a sharp predictor of $y^{\mathrm{base}}$.
Evaluated on coin-flip labels its Brier rises by $0.084$ on average, while
the same forest evaluated on a linear tail that covers $23\%$ of the window
rises by only $0.033$.
Size matching does not by itself fix the comparison.
On the region that attains $\Delta^{\mathrm{cf}}$, the size-matched excess
on Income noise remains $0.071$, against $0.002$ on clean sites.
Multiplying by the skill factor is what sends the noise score from $0.071$
to $0.001$.
Public Coverage does not reproduce the $0.40$ figure.
There the baseline concept is harder, the pooled excess on noise is
$0.035$, and every structured site still outranks it.
The failure mode is real on the task where the reference fit is strong,
which is the task on which an ungated risk gap is most tempting.

\begin{figure}[t]
\centering
\includegraphics[width=\textwidth]{concordance_cd_vs_clean.png}
\caption{Pairwise concordance against clean sites.
The dashed line is the $0.8$ target.
$\Delta^{\mathrm{cf}}$, the whole-window score, skill-gated $\tau^2$, raw
$\tau^2$, and the Pareto rank clear the line on both tasks.
The joint classifier two-sample test also clears it on average; its failure
is the noise contrast in Table~\ref{tab:concordance}.
Label total variation, the random-forest domain AUC, and RBF MMD do not
clear the line.
Source: abrupt activation, $\alpha=1$, $q=0.5$, four seeds.}
\label{fig:concordance}
\end{figure}

The Pareto comparison uses three axes when the time curve below is
available, and two otherwise.
The axes are $\Delta^{\mathrm{cf}}$, $\tau^{2,\mathrm{g}}$, and the time
jump, defined as the maximum of the curve on endpoints at or after
$t^\star$ minus the maximum on endpoints before $t^\star$.
A site's rank is one plus the number of sites that are at least as large
on every axis and strictly larger on at least one.
Concordance is computed on the negative rank, so a smaller rank is a
higher score.
Mean ranks on Income are $1.875$ for nonlinear sites, $3.250$ for linear
sites, $6.625$ for clean sites, and $6.750$ for the noise site.
On Public Coverage they are $1.375$, $2.667$, $7.188$, and $6.500$.
The undominated front is mostly the nonlinear sites: half of them on
Income and five eighths on Public Coverage.
No clean site and no noise site is undominated.
Linear sites are often dominated by a nonlinear site, because the
nonlinear rewrite covers every row and the linear rewrite covers a tail,
but they remain strictly ahead of every clean site.
That is the sense in which the Pareto order puts primary sources first.
It does not claim that every primary source lies on the front.

\subsection{Time localization}
\label{sec:time}

The curve at site $k$ reuses the burn-in reference and recomputes
\eqref{eq:delta-cf} on a sliding window of four batches, using only the
whole window and the upper tail.
The horizontal coordinate in Figure~\ref{fig:time} is the last batch
included in the window.
Drift begins at batch $8$.
A window of length four is entirely inside the drift only when its last
batch is at least $11$.
The figure follows that arithmetic.
Means for clean sites and for the noise site stay below $0.01$ on both
sides of $t^\star$.
Linear and nonlinear means leave zero between batches $9$ and $10$ and
level off at batch $11$.
On Income the nonlinear plateau is near $0.32$ and the linear plateau is
near $0.22$.
On Public Coverage the two plateaus sit together near $0.18$, which is
consistent with the similar headline scores in Table~\ref{tab:means}.

\begin{figure}[t]
\centering
\includegraphics[width=\textwidth]{time_profile_delta_cf.png}
\caption{Mean $\Delta^{\mathrm{cf}}$ against the last batch in a window of
length four.
The vertical line is $t^\star=8$.
The curve levels off at batch $11$, the first endpoint whose window lies
entirely after the change.
Clean and noise traces remain near zero.
Source: abrupt activation, $\alpha=1$, $q=0.5$, four seeds.}
\label{fig:time}
\end{figure}

Two detectors are applied to each curve without being told $t^\star$.
The online rule calibrates on the first two evaluated endpoints, sets the
threshold at that maximum plus $0.01$, and reports the first later endpoint
that exceeds the threshold.
A hit is an estimate within two batches of $t^\star$.
The derivative rule reports the endpoint of the largest first difference.
Table~\ref{tab:time} shows that the derivative is estimating the onset of
the plateau rather than the onset of the drift.
Its mean absolute error is about three batches on every structured type,
which is one window length.
The online rule is the onset estimator.
It hits every nonlinear site on both tasks and every linear site on Public
Coverage.
Income linear sites, whose tail is the thin one, hit in $0.833$ of
replicates, with mean absolute error $1.92$ batches.
The cost is a false alarm rate of $0.375$ on Income clean sites and
$0.125$ on Public Coverage clean sites: any crossing counts, whether or
not it falls near $t^\star$.
Among the Income clean curves that do cross, the mean absolute error is
$4.17$ batches, so the typical alarm is not a call at the change.
A threshold of $0.01$ above a two-point calibration is enough to rank
sites when the window is the full recent block, and it is not enough to
call a changepoint on a short window without a second screen.
The noise site crosses the threshold in two of four Income replicates and
in one of four Public Coverage replicates.
One Income crossing lands on batch $10$, inside the two-batch window; the
other is six batches later.
The Public Coverage crossing is three batches after $t^\star$ and is not
counted as a hit.

\begin{table}[t]
\centering
\caption{Time localization over four seeds.
Online hit rate is the fraction of sites whose first threshold crossing
lies within two batches of $t^\star$.
False-alarm rate is computed only for clean sites.
Derivative MAE is the absolute error of the largest first difference.}
\label{tab:time}
\begin{tabular}{llcccc}
\toprule
Task & Type & Online hit & False alarm & Online MAE & Derivative MAE \\
\midrule
Income & Clean & 0.062 & 0.375 & 4.17 & 4.06 \\
Income & Linear & 0.833 & --- & 1.92 & 2.83 \\
Income & Nonlinear & 1.000 & --- & 1.62 & 3.00 \\
Income & Noise & 0.250 & --- & 4.00 & 3.25 \\
Public Coverage & Clean & 0.062 & 0.125 & 3.00 & 3.88 \\
Public Coverage & Linear & 1.000 & --- & 1.42 & 3.00 \\
Public Coverage & Nonlinear & 1.000 & --- & 1.50 & 3.00 \\
Public Coverage & Noise & 0.000 & --- & 3.00 & 4.00 \\
\bottomrule
\end{tabular}
\end{table}

\subsection{Concept-group type resolution}
\label{sec:type}

Detecting a primary source does not say whether the new concept is the
linear halfspace or one of the two nonlinear rules.
The dictionary is fixed before looking at site labels.
On a window with outcome $y$, three logistic regressions are cross-fit
against an intercept-only Brier score:
$(Z_1,Z_2)$, the scalar $Z_1 Z_2$, and $(\sin Z_1,\, Z_3^2)$.
Group importance is the drop in Brier.
The call is the maximizing group.
One call uses every row in the recent window.
The other uses only the upper tail of $X_1$, the region in which the
linear rule was planted.
Accuracy is computed on the five structured sites only.
Top-five precision asks a different question: of the five sites with the
largest $\Delta^{\mathrm{cf}}$, how many lie in $\mathcal{P}$.

\begin{table}[t]
\centering
\caption{Type resolution at $\alpha=1$, abrupt activation, and $q=0.5$.
Accuracy is the fraction of structured sites whose maximizing concept group
equals the planted mechanism, averaged over four seeds.
Recall is pooled over both tasks and all four seeds, hence over $24$ linear
sites and $8$ sites of each nonlinear mechanism.
Top-five precision of $\Delta^{\mathrm{cf}}$ is $1$ in every seed on both
tasks.}
\label{tab:type}
\begin{tabular}{llcc}
\toprule
Call & Mechanism or task & Accuracy or recall \\
\midrule
All rows & Income, accuracy & 0.700 \\
All rows & Public Coverage, accuracy & 0.700 \\
Upper tail & Income, accuracy & 0.750 \\
Upper tail & Public Coverage, accuracy & 0.900 \\
\midrule
All rows & Linear, recall & 0.500 \\
All rows & Product, recall & 1.000 \\
All rows & Sine-square, recall & 1.000 \\
Upper tail & Linear, recall & 0.917 \\
Upper tail & Product, recall & 0.375 \\
Upper tail & Sine-square, recall & 1.000 \\
\bottomrule
\end{tabular}
\end{table}

Table~\ref{tab:type} separates detection from typing.
Top-five precision is $1$ in every seed: the five largest values of
$\Delta^{\mathrm{cf}}$ are exactly the five structured sites.
Conditional on that detection, the full-window dictionary recovers both
nonlinear mechanisms in every replicate and recovers the linear halfspace
in half of the linear sites.
The miss is structural.
Outside the tail the linear site still carries $y^{\mathrm{base}}$, and on
Income that baseline outcome is already monotone in hours and schooling.
An additive logistic fit can look competitive with the product and
sine-square groups even when the rewritten slice is a halfspace in
$(Z_1,Z_2)$.
Restricting the call to the upper tail raises linear recall from $0.500$
to $0.917$ and raises Public Coverage accuracy from $0.700$ to $0.900$.
The same restriction lowers product recall from $1$ to $0.375$.
The product rule was planted on every row.
Inside the upper tail the factors are truncated, and the additive group
in $(Z_1,Z_2)$ is no longer dominated by the product scalar.
The sine-square group remains separated on either window.
Type resolution is therefore a property of the region on which the group
is scored, not a single accuracy number.

\subsection{Feature support and one-shot aggregation}
\label{sec:features}

Feature scores are computed on the recent window of each structured site
and compared with that site's planted support by AUC.
Four scores are used.
Shift permutation importance is the drop in the squared concept contrast
when a coordinate is permuted inside five quantile bins of the covariate
it is most correlated with.
Both the local forest and the reference forest are re-evaluated on that
same permuted matrix.
The local forest here is fit on the full recent window; the cross-fit in
\eqref{eq:delta-cf} is not reused.
Predictive permutation importance is the drop in held-out AUC for
predicting the recent label.
The difference of that importance between the recent window and the burn-in
is the fourth score's companion; the fourth score itself is the absolute coefficient vector of an
$\ell_1$-penalized logistic regression on the recent window, with
$C=0.5$.

\begin{table}[t]
\centering
\caption{Mean AUC of a feature score against the planted support.
Linear and product supports are $\{X_1,X_2\}$.
The sine-square support is $\{X_1,X_3\}$.
Averages run over the structured sites of four seeds.}
\label{tab:feature}
\begin{tabular}{llccc}
\toprule
Task & Type & Shift permutation & Predictive permutation & $\ell_1$ coefficient \\
\midrule
Income & Linear & 0.906 & 0.927 & 0.786 \\
Income & Nonlinear & 0.789 & 1.000 & 0.930 \\
Public Coverage & Linear & 0.953 & 0.995 & 0.841 \\
Public Coverage & Nonlinear & 0.985 & 1.000 & 0.879 \\
\bottomrule
\end{tabular}
\end{table}

Table~\ref{tab:feature} says that the planted coordinates are recoverable
from a single drifted site.
Predictive permutation AUC is at least $0.927$.
Shift permutation AUC is at least $0.906$ on linear sites and falls to
$0.789$ on Income nonlinear sites, where the product and the sine-square
rule are functions of coordinates that also predict the baseline income
label.
The $\ell_1$ coefficient is weaker than the forest on linear Public
Coverage ($0.841$) and stronger than shift permutation on Income nonlinear
sites ($0.930$).
None of these AUCs is a site-ranking result.
They are the contents of the vector a site would upload.

Aggregation treats each upload as a point on the simplex.
Negative shift-permutation coordinates are truncated at zero and the vector
is renormalized.
The target vector puts mass on a coordinate proportional to the number of
structured sites that plant it.
Uniform averaging, the coordinate-wise median, inverse-variance weights
with variance estimated from the permutation repeats and floored at
$10^{-4}$, and weights proportional to $\Delta^{\mathrm{cf}}$ are compared
by AUC, precision at the number of positive target coordinates, and
Spearman correlation with the target.
Leaving the noise site out of the uniform average, and averaging only the
five structured sites, are oracles used as ceilings rather than as methods.

On Public Coverage every shift aggregator has AUC at least $0.984$, and
the $\Delta^{\mathrm{cf}}$-weighted average and the average that drops the
noise site both have AUC $1$ and precision $1$.
On Income the uniform shift average has AUC $0.869$ and precision $0.667$.
Weighting by $\Delta^{\mathrm{cf}}$ raises those figures to $0.917$ and
$0.750$.
The predictive permutation upload and the $\ell_1$ upload, averaged
uniformly, reach AUC $0.988$ on Income.
That is not evidence against the concept score.
On Income the planted coordinates are hours, schooling, and age, which are
also the coordinates that predict $y^{\mathrm{base}}$ on a clean site.
A clean site's predictive vector already points at the target.
The shift vector is the one that is trying to upload the change rather than
the baseline regression, and it is the vector that improves when clean and
noise sites are down-weighted by $\Delta^{\mathrm{cf}}$.

Leave-one-site-out influence, defined as one minus the Spearman
correlation between the uniform aggregate and the aggregate with one site
removed, does not identify primary sources.
Mean influence on Income is $0.033$ for clean sites, $0.013$ for linear
sites, $0.018$ for the noise site, and $0.021$ for nonlinear sites.
On Public Coverage every type mean lies between $0.005$ and $0.008$.
The uploaded vectors agree with each other closely enough that deleting one
site barely moves the ranking.
Site-level $\Delta^{\mathrm{cf}}$ is the functional that separates primary
sources from clean sites.
Influence on a simplex average is not.

\subsection{Intensity, tail mass, and gradual activation}
\label{sec:sensitivity}

The main grid holds $\alpha=1$, abrupt activation, and $q=0.5$.
The sensitivity grid varies those three levers on Income only, with three
seeds and the same score, and it does not recompute feature uploads or the
joint two-sample test.
Because the $0.5$ and $0.7$ quantiles of hours worked are the same mass
point at $40$ hours, $q=0.7$ selects the same tail as $q=0.5$ and is
omitted.
The quantile $0.8$ moves the threshold to $45$ hours and leaves about
$17\%$ of rows in the linear tail.
The quantile $0.9$ moves it to $50$ hours and leaves about $8\%$.
Nonlinear sites continue to rewrite every activated row.
Only the linear sites become thinner.

\begin{table}[t]
\centering
\caption{Income concordance $C(\Delta^{\mathrm{cf}};\mathcal{P},\mathcal{C})$
under weaker rewrites and thinner linear tails.
Three seeds.
The worst seed is in parentheses.
Quantile $0.5$ is the main tail, about $23\%$ of rows because of the mass
at $40$ hours.}
\label{tab:sensitivity}
\begin{tabular}{lcccccc}
\toprule
& \multicolumn{2}{c}{$q=0.5$} & \multicolumn{2}{c}{$q=0.8$} & \multicolumn{2}{c}{$q=0.9$} \\
\cmidrule(lr){2-3}\cmidrule(lr){4-5}\cmidrule(lr){6-7}
$\alpha$ & Abrupt & Gradual & Abrupt & Gradual & Abrupt & Gradual \\
\midrule
1.0 & 1.000 {\scriptsize(1.000)} & 1.000 {\scriptsize(1.000)} & 1.000 {\scriptsize(1.000)} & 1.000 {\scriptsize(1.000)} & 0.850 {\scriptsize(0.750)} & 0.867 {\scriptsize(0.800)} \\
0.7 & 1.000 {\scriptsize(1.000)} & 1.000 {\scriptsize(1.000)} & 0.967 {\scriptsize(0.950)} & 0.950 {\scriptsize(0.900)} & 0.817 {\scriptsize(0.750)} & 0.683 {\scriptsize(0.500)} \\
0.4 & 0.783 {\scriptsize(0.650)} & 0.750 {\scriptsize(0.650)} & 0.617 {\scriptsize(0.450)} & 0.600 {\scriptsize(0.450)} & 0.550 {\scriptsize(0.450)} & 0.550 {\scriptsize(0.450)} \\
\bottomrule
\end{tabular}
\end{table}

\begin{figure}[t]
\centering
\includegraphics[width=\textwidth]{sensitivity_concordance.png}
\caption{Income concordance of $\Delta^{\mathrm{cf}}$ against clean sites.
Each panel is one linear-tail quantile.
Cell entries are means over three seeds.
Means stay at or above $0.8$ for $\alpha\in\{0.7,1\}$ down to a linear tail
of about $17\%$ of rows.
They fall below $0.8$ when $\alpha=0.4$, and when a tail of about $8\%$ is
combined with $\alpha=0.7$ and a gradual ramp.
Full intensity on that thin tail stays above $0.8$ on average and not in
every seed; the worst seeds are in Table~\ref{tab:sensitivity}.}
\label{fig:sensitivity}
\end{figure}

Table~\ref{tab:sensitivity} and Figure~\ref{fig:sensitivity} locate the
boundary of the $0.8$ target.
At full intensity the concordance remains $1$ down to a linear tail of
about $17\%$ of rows, whether the change is abrupt or spread over four
batches.
At a tail of about $8\%$, full abrupt intensity falls to mean concordance
$0.850$ with a worst seed of $0.750$, so the mean still clears $0.8$ and
the guarantee across seeds does not.
Gradual activation at that thin tail has the same character
($0.867$, worst seed $0.800$).
Lowering the rewrite probability to $0.7$ preserves the target on the main
tail and on the $17\%$ tail.
Combining $\alpha=0.7$ with the $8\%$ tail and a four-batch ramp drops the
concordance to $0.683$, with a worst seed of $0.500$.
At $\alpha=0.4$ every column is below $0.8$.
The main-tail abrupt entry is $0.783$ with a worst seed of $0.650$.
On that same $\alpha=0.4$ slice, the whole-window score without the tail
scan has concordance $0.433$.
The region maximum is what keeps a weak but structured rewrite near the
target.
It does not keep an arbitrarily weak rewrite above the target.

\subsection{What the benchmark supports}
\label{sec:support}

The benchmark supports four claims, each tied to a table.
First, when the rewrite is complete and the linear tail contains about a
fifth of the rows or more, $\Delta^{\mathrm{cf}}$ orders every structured
site ahead of every clean site and ahead of the noise site, on two real
covariate tables, in every seed
(Table~\ref{tab:concordance}).
Second, the functionals that ignore either the existence of a local concept
or the region in which that concept lives do not.
Label total variation and the covariate arms are near chance against clean
sites.
The joint two-sample test and the pooled Brier excess promote the noise
site, completely so for linear sites under the pooled excess on Income.
Third, the time profile of the same score locates an abrupt change to
within about two batches once the scoring window is taken into account,
with a false-alarm rate on clean Income sites that a deployment threshold
would still have to budget for
(Table~\ref{tab:time}).
Fourth, a pre-specified concept-group dictionary can name the nonlinear
mechanisms from the full window and the linear mechanism from the tail, and
it cannot do both from a single region
(Table~\ref{tab:type}).

The benchmark does not support a claim that leave-one-site-out influence
on an averaged variable-importance vector identifies the primary site.
Those vectors are too similar across sites for the deletion of one upload
to move the rank correlation.
It also does not support a claim at very small intensity.
Below a rewrite probability of one half, or inside a linear tail that
covers less than a tenth of the window, pairwise concordance falls through
$0.8$ on at least one seed.
Both limits are visible in the same design that produces the perfect
separation at full intensity, which is the reason for reporting them beside
the headline rather than as a separate simulation.

% after defining nothing: the section is self-contained.

\section{Primary-source identification across heterogeneous sites}
\label{sec:primary-source}

The manuscript separates covariate shift, a change in $P(X)$, from concept
drift, a change in $P(Y\mid X)$ between a reference batch
$D_{\mathrm{exist}}$ and a new batch $D_{\mathrm{new}}$.
Feature-level procedures in the paper, including the pseudo-outcome
risk and conditional permutation importance, score coordinates inside one
batch.
The federated aggregation in Appendix~11 then asks whether those
variable-importance vectors can be merged without moving raw rows.
That merge treats every site as a noisy look at one common shift.
A prior operational question is which sites are the primary sources of the
shift, and which sites are either stable or actively destroying the
consensus.
This section constructs a ten-site benchmark in which the answer is known,
and evaluates a counterfactual concept score against that answer.

The design isolates concept drift.
Every site is a random slice of one real covariate table, so the marginal
distribution $P(X)$ is shared.
The American Community Survey extracts below are cross-sections, not
panels, so the batch index is a seeded random order rather than a clock.
Stationarity of $P(X)$ is then a construction, and a covariate arm that
fires is evidence of a bug in the generator rather than a feature of the
task.
Labels follow the real outcome through a burn-in and are then rewritten,
on a known subset of sites, by a known rule.

\subsection{Data and the ten-site generator}
\label{sec:dgp}

Two WhyShift tasks supply the base table $(X,y^{\mathrm{base}})$, both
taken from the 2018 ACS one-year person file for California and both using
the official \texttt{folktables} feature lists.
ACS Income has $p=10$ covariates and the indicator that personal income
exceeds $\$50{,}000$.
ACS Public Coverage has $p=19$ covariates and the indicator of public
health-insurance coverage.
In each Monte Carlo replicate the analysis sample is winsorized at the
first and ninety-ninth percentiles, standardized, and cut into $K=10$
non-overlapping sites of $n_k=1600$ rows.
Each site is then cut into $T=16$ batches of $100$ rows.
The burn-in is the first four batches.
The concept change, when it is switched on, begins at batch
$t^\star=8$.

The ten sites have four roles.
Four clean sites keep $y=y^{\mathrm{base}}$ at every batch.
Three linear sites rewrite labels inside an upper tail of a designated
coordinate $X_1$.
Two nonlinear sites rewrite every activated row.
One noise site replaces activated labels with independent
$\mathrm{Bernoulli}(1/2)$ draws.
The noise site is a consensus breaker.
It is not a primary source: there is no new concept for a local model to
learn.
Table~\ref{tab:roles} records which coordinates carry the rules.
Income uses hours worked, educational attainment, and age.
Public Coverage uses age, educational attainment, and personal income.
These are the continuous or ordinal columns in each official feature list.
The remaining columns stay in the table as nuisance covariates.

\begin{table}[t]
\centering
\caption{Rule coordinates. $Z_j$ denotes the analysis-sample standardize
of the named column. Linear sites rewrite only the upper tail of $Z_1$.
Nonlinear sites rewrite every activated row.}
\label{tab:roles}
\begin{tabular}{llccc}
\toprule
Task & Role of $(X_1,X_2,X_3)$ & $X_1$ & $X_2$ & $X_3$ \\
\midrule
ACS Income & hours, schooling, age & WKHP & SCHL & AGEP \\
ACS Public Coverage & age, schooling, income & AGEP & SCHL & PINCP \\
\bottomrule
\end{tabular}
\end{table}

Write $Z$ for the standardized covariates and let
$M_k(q)=\{Z_1>Q_q(Z_1)\}$ be the upper tail at quantile $q$, computed on
the pooled analysis sample.
The main design uses $q=0.5$.
On a linear site the candidate label is the halfspace
\begin{equation}
\label{eq:linear-rule}
y^{\mathrm{rule}}
=
\mathbf{1}\!\left\{
\sigma_k(\beta_{k1} Z_1+\beta_{k2} Z_2)>c_k
\right\}
\quad\text{on }M_k(q),
\end{equation}
and $y^{\mathrm{base}}$ off that tail.
The three coefficient vectors are
$(1,0.35)$, $(0.25,1)$, and $(-0.85,0.55)$.
The threshold $c_k$ is the quantile of the linear score that matches the
positive rate of $y^{\mathrm{base}}$ on $M_k(q)$, and the sign
$\sigma_k\in\{\pm 1\}$ is the orientation whose matched labels agree less
with $y^{\mathrm{base}}$.
The sign choice is part of the generator.
It keeps the planted concept from collapsing back onto the real outcome
when hours or schooling already predict income.
It is not estimated by the detector.

On a nonlinear site the same rate match is applied to every row, not only
the tail:
\begin{equation}
\label{eq:nonlinear-rule}
y^{\mathrm{rule}}
=
\mathbf{1}\{Z_1 Z_2>c\}
\qquad\text{or}\qquad
y^{\mathrm{rule}}
=
\mathbf{1}\{\sin(Z_1)+Z_3^2>c\}.
\end{equation}
The first site uses the product and the second uses the sine-square rule.
The planted feature support is $\{X_1,X_2\}$ for every linear site and for
the product site, and $\{X_1,X_3\}$ for the sine-square site.

Activation is a batch weight $w_t\in[0,1]$.
Under abrupt activation, $w_t=\mathbf{1}\{t\ge t^\star\}$.
Under gradual activation, $w_t=0$ before $t^\star$ and then rises linearly
to one over four batches.
Each post-burn-in row is rewritten with probability $\alpha w_t$, and
otherwise keeps $y^{\mathrm{base}}$.
The main grid uses $\alpha=1$.
The noise site uses the same weights and then draws a fresh coin flip.
Clean sites ignore the weight.

Two consequences are worth stating before any score is computed.
First, because the structured rules match the local positive rate, a
marginal label discrepancy is not a substitute for concept drift.
The largest average gap between the post-change positive rate and the
baseline, among structured sites, is $2.4$ percentage points
(Income, nonlinear).
Second, hours worked in the ACS piles up at $40$.
The empirical quantiles of WKHP at $0.5$ and at $0.7$ are both $40$, so
$\{Z_1>Q_{0.5}\}$ and $\{Z_1>Q_{0.7}\}$ are the same set, about $23\%$ of
Income rows.
The median tail in the main Income design is therefore thinner than a
continuous coordinate would suggest.
Public Coverage uses age as $X_1$, and that tail is about half the sample.
Table~\ref{tab:dgp-check} records the realized flip rates.
Structured sites move roughly half the labels inside the region they
rewrite, and they leave the marginal positive rate near the baseline.
The noise site flips about half of all post-burn-in labels and shifts the
positive rate to one half.

\begin{table}[t]
\centering
\caption{Realized label changes after $t^\star$, averaged over four seeds
at $\alpha=1$, abrupt activation, and $q=0.5$.
Flip rate is the fraction of post-burn-in rows whose label differs from
$y^{\mathrm{base}}$.
Agreement is the fraction of the rewritten region on which the rate-matched
rule reproduces $y^{\mathrm{base}}$.
Clean sites are not rewritten, so agreement is one.
The noise site has no rule, so agreement is left blank.}
\label{tab:dgp-check}
\begin{tabular}{llcccc}
\toprule
Task & Site type & Flip rate & Rule agreement & Post positive rate & Base positive rate \\
\midrule
Income & Clean & 0.000 & 1.000 & 0.406 & 0.406 \\
Income & Linear & 0.131 & 0.427 & 0.400 & 0.414 \\
Income & Nonlinear & 0.555 & 0.436 & 0.396 & 0.420 \\
Income & Noise & 0.498 & --- & 0.491 & 0.410 \\
Public Coverage & Clean & 0.000 & 1.000 & 0.369 & 0.367 \\
Public Coverage & Linear & 0.259 & 0.476 & 0.378 & 0.372 \\
Public Coverage & Nonlinear & 0.492 & 0.506 & 0.371 & 0.374 \\
Public Coverage & Noise & 0.496 & --- & 0.495 & 0.362 \\
\bottomrule
\end{tabular}
\end{table}

\subsection{The counterfactual concept score}
\label{sec:delta-cf}

Let $\mathcal{D}_0$ be the pooled burn-in, eight thousand rows in the main
design, still generated from $y^{\mathrm{base}}$ at every site.
Let $\mathcal{D}_k$ be the most recent six batches at site $k$, six hundred
rows, all of them inside the drift window when activation is abrupt.
A random forest with $25$ trees, depth $8$, and minimum leaf size $5$ is
the outcome model throughout, in line with the recommendation that the
forest be used for distance estimation rather than for a causal reading of
the meta-learner.
Local predictions on $\mathcal{D}_k$ are three-fold cross-fit probabilities
$\mu_k^{\mathrm{oof}}$.
The reference probabilities $\mu_0^{\mathrm{sm}}$ are an average of two
forests, each trained on a random subset of $\mathcal{D}_0$ whose size
equals the local training set, two thirds of the current window.
Matching the training size stops the larger reference sample from looking
like a concept change.
A third forest trained on all of $\mathcal{D}_0$ is retained only as the
pooled excess defined below.
It is not the headline score.

On a region $\mathcal{R}$ that has both classes in the current window and
in the reference, at least fifty current rows, and at least eighty
reference rows, write $B(y,p)$ for the Brier score and
\begin{equation}
\label{eq:delta-region}
\Delta_k(\mathcal{R})
=
\Big(B(\mu_0^{\mathrm{sm}})-B(\mu_k^{\mathrm{oof}})\Big)_+
\cdot
\Big(2\big(\mathrm{AUC}(\mu_k^{\mathrm{oof}})-{\textstyle\frac12}\big)\Big)_+.
\end{equation}
The first factor is the excess risk of the reference concept over a local
concept of the same capacity.
The second factor is zero when the local model does not beat a coin, and
it equals one when the local AUC equals one.
The headline score is the maximum over a dictionary that does not use
labels,
\begin{equation}
\label{eq:delta-cf}
\Delta_k^{\mathrm{cf}}
=
\max_{\mathcal{R}\in\{\mathrm{all},\,\mathrm{high},\,\mathrm{low}\}}
\Delta_k(\mathcal{R}),
\end{equation}
where ``high'' and ``low'' are the two sides of the burn-in median of
$X_1$.
The median is computed on $\mathcal{D}_0$ only.
The whole-window score $\Delta_k(\mathrm{all})$ is reported beside the
maximum so that the cost of ignoring the tail is visible.

The same two probability vectors define the concept contrast
\begin{equation}
\label{eq:tau}
\tau_k^2
=
\frac{1}{|\mathcal{D}_k\cap\mathcal{R}_k^\star|}
\sum_{i\in\mathcal{D}_k\cap\mathcal{R}_k^\star}
\big(\mu_k^{\mathrm{oof}}(X_i)-\mu_0^{\mathrm{sm}}(X_i)\big)^2,
\end{equation}
with $\mathcal{R}_k^\star$ the region attaining the maximum in
\eqref{eq:delta-cf}.
This is the concept contrast that the pseudo-outcome decomposition isolates
from the covariate arm once $P(X)$ is shared across $D_{\mathrm{exist}}$
and $D_{\mathrm{new}}$.
The benchmark does not re-estimate a propensity: under the common
covariate table that weight is not where the signal sits.
Gating by the same skill factor as in \eqref{eq:delta-region} produces
$\tau_k^{2,\mathrm{g}}$.
Raw $\tau_k^2$ is retained because a flat local predictor on a noise site
can still sit far from a sharp reference predictor.
The gate is the part of $\Delta_k^{\mathrm{cf}}$ that refuses to call that
gap a primary source.

For a score $s$ and two disjoint collections of sites $\mathcal{A}$ and
$\mathcal{B}$, the pairwise concordance is
\begin{equation}
\label{eq:concordance}
C(s;\mathcal{A},\mathcal{B})
=
\frac{1}{|\mathcal{A}|\,|\mathcal{B}|}
\sum_{a\in\mathcal{A}}\sum_{b\in\mathcal{B}}
\mathbf{1}\{s_a>s_b\}.
\end{equation}
Ties count as failures.
The primary set $\mathcal{P}$ is the five structured sites.
The target for the headline score is
$C(\Delta^{\mathrm{cf}};\mathcal{P},\mathcal{C})>0.8$,
where $\mathcal{C}$ is the set of four clean sites.
The same functional is reported for $\mathcal{P}$ against the noise site,
and separately for the linear and nonlinear subsets of $\mathcal{P}$.
Every entry below is a mean over four seeds, with the worst seed in
parentheses, unless a row is explicitly a per-type mean.

Negative controls use the same windows and never see the label rule.
The domain AUC is the cross-fit AUC of a classifier trained to separate
burn-in rows from recent rows using $X$ alone, once with the forest and
once with logistic regression.
The covariate discrepancy is an RBF MMD on subsamples of size $250$, with
bandwidth equal to the median of the squared Euclidean distances.
A joint classifier two-sample test appends $y$ to $X$ before the same
forest AUC is computed.
Marginal label total variation is the absolute difference of positive rates
between the recent window and the burn-in.

\begin{table}[t]
\centering
\caption{Mean score by site type at $\alpha=1$, abrupt activation, and
$q=0.5$. Four seeds. $\Delta^{\mathrm{cf}}$ is \eqref{eq:delta-cf}.
$\Delta(\mathrm{all})$ is the whole-window score.
$\Delta(\mathrm{high})$ is the upper tail of $X_1$.
Domain AUC is the logistic classifier on $X$ only.}
\label{tab:means}
\small
\begin{tabular}{llrrrrrr}
\toprule
Task & Type & $\Delta^{\mathrm{cf}}$ & $\Delta(\mathrm{all})$ & $\Delta(\mathrm{high})$ & $\tau^2$ & Skill & Domain AUC \\
\midrule
Income & Clean & 0.003 & 0.000 & 0.002 & 0.015 & 0.696 & 0.498 \\
Income & Linear & 0.253 & 0.024 & 0.253 & 0.181 & 0.937 & 0.482 \\
Income & Nonlinear & 0.336 & 0.298 & 0.316 & 0.313 & 0.897 & 0.494 \\
Income & Noise & 0.001 & 0.000 & 0.000 & 0.096 & 0.008 & 0.494 \\
Public Coverage & Clean & 0.001 & 0.000 & 0.001 & 0.014 & 0.364 & 0.503 \\
Public Coverage & Linear & 0.200 & 0.052 & 0.200 & 0.132 & 0.943 & 0.493 \\
Public Coverage & Nonlinear & 0.236 & 0.182 & 0.179 & 0.159 & 0.970 & 0.478 \\
Public Coverage & Noise & 0.003 & 0.002 & 0.003 & 0.046 & 0.089 & 0.528 \\
\bottomrule
\end{tabular}
\end{table}

Table~\ref{tab:means} is the raw material of the concordance claim.
On both tasks, and for every linear site, the maximum is exactly the
upper-tail score: $\Delta^{\mathrm{cf}}=\Delta(\mathrm{high})$.
On Income that tail score is $0.253$ against a whole-window score of
$0.024$.
The whole window is mostly unrewritten rows, and a single Brier average
dilutes a real halfspace by about an order of magnitude.
The noise skill is $0.008$ on Income and $0.089$ on Public Coverage, so
the gate shuts even though the ungated contrast $\tau^2$ on Income noise
($0.096$) exceeds the clean contrast ($0.015$).
Logistic domain AUCs sit on $1/2$ for every type.
The largest is $0.528$, on the Public Coverage noise site.
The forest domain AUC, not shown column by column, has type means between
$0.490$ and $0.514$, and equals $0.514$ on that same noise site.
Neither is evidence of a covariate shift.

\subsection{Concordance}
\label{sec:concordance}

\begin{table}[t]
\centering
\caption{Pairwise concordance \eqref{eq:concordance} at $\alpha=1$, abrupt
activation, and $q=0.5$.
Each cell is the mean over four seeds, with the worst seed in parentheses.
$\mathcal{P}$ is the five structured sites and $\mathcal{C}$ is the four
clean sites.}
\label{tab:concordance}
\small
\begin{tabular}{llcccc}
\toprule
Task & Score & $\mathcal{P}>\mathcal{C}$ & $\mathcal{P}>\mathrm{noise}$ & Linear $>\mathcal{C}$ & Nonlinear $>\mathcal{C}$ \\
\midrule
Income & $\Delta^{\mathrm{cf}}$ & 1.000 {\scriptsize(1.000)} & 1.000 {\scriptsize(1.000)} & 1.000 {\scriptsize(1.000)} & 1.000 {\scriptsize(1.000)} \\
Income & $\Delta(\mathrm{all})$ & 1.000 {\scriptsize(1.000)} & 1.000 {\scriptsize(1.000)} & 1.000 {\scriptsize(1.000)} & 1.000 {\scriptsize(1.000)} \\
Income & $\tau^{2,\mathrm{g}}$ & 1.000 {\scriptsize(1.000)} & 1.000 {\scriptsize(1.000)} & 1.000 {\scriptsize(1.000)} & 1.000 {\scriptsize(1.000)} \\
Income & Pareto rank & 1.000 {\scriptsize(1.000)} & 1.000 {\scriptsize(1.000)} & 1.000 {\scriptsize(1.000)} & 1.000 {\scriptsize(1.000)} \\
Income & $\tau^2$ & 1.000 {\scriptsize(1.000)} & 1.000 {\scriptsize(1.000)} & 1.000 {\scriptsize(1.000)} & 1.000 {\scriptsize(1.000)} \\
Income & Pooled excess & 1.000 {\scriptsize(1.000)} & 0.400 {\scriptsize(0.400)} & 1.000 {\scriptsize(1.000)} & 1.000 {\scriptsize(1.000)} \\
Income & Label TV & 0.500 {\scriptsize(0.100)} & 0.050 {\scriptsize(0.000)} & 0.500 {\scriptsize(0.083)} & 0.500 {\scriptsize(0.125)} \\
Income & Joint C2ST & 0.850 {\scriptsize(0.800)} & 0.350 {\scriptsize(0.200)} & 0.750 {\scriptsize(0.667)} & 1.000 {\scriptsize(1.000)} \\
Income & RF domain AUC & 0.487 {\scriptsize(0.400)} & 0.500 {\scriptsize(0.000)} & 0.521 {\scriptsize(0.333)} & 0.438 {\scriptsize(0.250)} \\
Income & RBF MMD & 0.412 {\scriptsize(0.300)} & 0.500 {\scriptsize(0.000)} & 0.375 {\scriptsize(0.167)} & 0.469 {\scriptsize(0.125)} \\
\midrule
Public Coverage & $\Delta^{\mathrm{cf}}$ & 1.000 {\scriptsize(1.000)} & 1.000 {\scriptsize(1.000)} & 1.000 {\scriptsize(1.000)} & 1.000 {\scriptsize(1.000)} \\
Public Coverage & $\Delta(\mathrm{all})$ & 1.000 {\scriptsize(1.000)} & 1.000 {\scriptsize(1.000)} & 1.000 {\scriptsize(1.000)} & 1.000 {\scriptsize(1.000)} \\
Public Coverage & $\tau^{2,\mathrm{g}}$ & 1.000 {\scriptsize(1.000)} & 1.000 {\scriptsize(1.000)} & 1.000 {\scriptsize(1.000)} & 1.000 {\scriptsize(1.000)} \\
Public Coverage & Pareto rank & 1.000 {\scriptsize(1.000)} & 1.000 {\scriptsize(1.000)} & 1.000 {\scriptsize(1.000)} & 1.000 {\scriptsize(1.000)} \\
Public Coverage & $\tau^2$ & 1.000 {\scriptsize(1.000)} & 1.000 {\scriptsize(1.000)} & 1.000 {\scriptsize(1.000)} & 1.000 {\scriptsize(1.000)} \\
Public Coverage & Pooled excess & 1.000 {\scriptsize(1.000)} & 1.000 {\scriptsize(1.000)} & 1.000 {\scriptsize(1.000)} & 1.000 {\scriptsize(1.000)} \\
Public Coverage & Label TV & 0.613 {\scriptsize(0.350)} & 0.000 {\scriptsize(0.000)} & 0.604 {\scriptsize(0.417)} & 0.625 {\scriptsize(0.250)} \\
Public Coverage & Joint C2ST & 0.950 {\scriptsize(0.850)} & 0.450 {\scriptsize(0.200)} & 0.917 {\scriptsize(0.750)} & 1.000 {\scriptsize(1.000)} \\
Public Coverage & RF domain AUC & 0.438 {\scriptsize(0.350)} & 0.200 {\scriptsize(0.000)} & 0.396 {\scriptsize(0.167)} & 0.500 {\scriptsize(0.250)} \\
Public Coverage & RBF MMD & 0.600 {\scriptsize(0.500)} & 0.200 {\scriptsize(0.000)} & 0.562 {\scriptsize(0.417)} & 0.656 {\scriptsize(0.375)} \\
\bottomrule
\end{tabular}
\end{table}

Table~\ref{tab:concordance} is the primary result.
On both tasks and in every seed,
$C(\Delta^{\mathrm{cf}};\mathcal{P},\mathcal{C})=1$,
and the same equality holds for linear sites alone, for nonlinear sites
alone, and for $\mathcal{P}$ against the noise site.
The whole-window score and the skill-gated contrast $\tau^{2,\mathrm{g}}$
match this ranking at full intensity.
Figure~\ref{fig:concordance} plots the clean-site comparison.
The covariate rows do not.
Random-forest domain AUC and RBF MMD have concordances between $0.38$ and
$0.66$, with worst seeds as low as $0.17$.
Label total variation is consistent with a matched positive rate: its
concordance against clean sites is $0.50$ on Income and $0.61$ on Public
Coverage, and its concordance against the noise site is $0.05$ and $0.00$.
The noise site is the unique marginal shift, and a marginal functional
ranks it first.

The joint classifier two-sample test sees any change in $(X,Y)$.
It separates $\mathcal{P}$ from $\mathcal{C}$ at $0.85$ on Income and $0.95$
on Public Coverage, which clears $0.8$ on average, but it ranks the noise
site as the stronger departure: concordance of $\mathcal{P}$ against noise
is $0.35$ and $0.45$.
On Income the linear-only comparison for the joint test falls to $0.75$.
A test of the joint distribution is not a test for a learnable new concept.

Pooled excess makes the same mistake in a sharper way, and only on the
task where the reference concept is easy.
It is the Brier gap between the forest trained on the entire burn-in and
the local cross-fit model, before size matching and before the skill gate.
On Income its concordance against the noise site is $0.400$ in every seed.
That fraction is exact: both nonlinear sites beat the noise site, and none
of the three linear sites do.
The pooled Income forest is a sharp predictor of $y^{\mathrm{base}}$.
Evaluated on coin-flip labels its Brier rises by $0.084$ on average, while
the same forest evaluated on a linear tail that covers $23\%$ of the window
rises by only $0.033$.
Size matching does not by itself fix the comparison.
On the region that attains $\Delta^{\mathrm{cf}}$, the size-matched excess
on Income noise remains $0.071$, against $0.002$ on clean sites.
Multiplying by the skill factor is what sends the noise score from $0.071$
to $0.001$.
Public Coverage does not reproduce the $0.40$ figure.
There the baseline concept is harder, the pooled excess on noise is
$0.035$, and every structured site still outranks it.
The failure mode is real on the task where the reference fit is strong,
which is the task on which an ungated risk gap is most tempting.

\begin{figure}[t]
\centering
\includegraphics[width=\textwidth]{concordance_cd_vs_clean.png}
\caption{Pairwise concordance against clean sites.
The dashed line is the $0.8$ target.
$\Delta^{\mathrm{cf}}$, the whole-window score, skill-gated $\tau^2$, raw
$\tau^2$, and the Pareto rank clear the line on both tasks.
The joint classifier two-sample test also clears it on average; its failure
is the noise contrast in Table~\ref{tab:concordance}.
Label total variation, the random-forest domain AUC, and RBF MMD do not
clear the line.
Source: abrupt activation, $\alpha=1$, $q=0.5$, four seeds.}
\label{fig:concordance}
\end{figure}

The Pareto comparison uses three axes when the time curve below is
available, and two otherwise.
The axes are $\Delta^{\mathrm{cf}}$, $\tau^{2,\mathrm{g}}$, and the time
jump, defined as the maximum of the curve on endpoints at or after
$t^\star$ minus the maximum on endpoints before $t^\star$.
A site's rank is one plus the number of sites that are at least as large
on every axis and strictly larger on at least one.
Concordance is computed on the negative rank, so a smaller rank is a
higher score.
Mean ranks on Income are $1.875$ for nonlinear sites, $3.250$ for linear
sites, $6.625$ for clean sites, and $6.750$ for the noise site.
On Public Coverage they are $1.375$, $2.667$, $7.188$, and $6.500$.
The undominated front is mostly the nonlinear sites: half of them on
Income and five eighths on Public Coverage.
No clean site and no noise site is undominated.
Linear sites are often dominated by a nonlinear site, because the
nonlinear rewrite covers every row and the linear rewrite covers a tail,
but they remain strictly ahead of every clean site.
That is the sense in which the Pareto order puts primary sources first.
It does not claim that every primary source lies on the front.

\subsection{Time localization}
\label{sec:time}

The curve at site $k$ reuses the burn-in reference and recomputes
\eqref{eq:delta-cf} on a sliding window of four batches, using only the
whole window and the upper tail.
The horizontal coordinate in Figure~\ref{fig:time} is the last batch
included in the window.
Drift begins at batch $8$.
A window of length four is entirely inside the drift only when its last
batch is at least $11$.
The figure follows that arithmetic.
Means for clean sites and for the noise site stay below $0.01$ on both
sides of $t^\star$.
Linear and nonlinear means leave zero between batches $9$ and $10$ and
level off at batch $11$.
On Income the nonlinear plateau is near $0.32$ and the linear plateau is
near $0.22$.
On Public Coverage the two plateaus sit together near $0.18$, which is
consistent with the similar headline scores in Table~\ref{tab:means}.

\begin{figure}[t]
\centering
\includegraphics[width=\textwidth]{time_profile_delta_cf.png}
\caption{Mean $\Delta^{\mathrm{cf}}$ against the last batch in a window of
length four.
The vertical line is $t^\star=8$.
The curve levels off at batch $11$, the first endpoint whose window lies
entirely after the change.
Clean and noise traces remain near zero.
Source: abrupt activation, $\alpha=1$, $q=0.5$, four seeds.}
\label{fig:time}
\end{figure}

Two detectors are applied to each curve without being told $t^\star$.
The online rule calibrates on the first two evaluated endpoints, sets the
threshold at that maximum plus $0.01$, and reports the first later endpoint
that exceeds the threshold.
A hit is an estimate within two batches of $t^\star$.
The derivative rule reports the endpoint of the largest first difference.
Table~\ref{tab:time} shows that the derivative is estimating the onset of
the plateau rather than the onset of the drift.
Its mean absolute error is about three batches on every structured type,
which is one window length.
The online rule is the onset estimator.
It hits every nonlinear site on both tasks and every linear site on Public
Coverage.
Income linear sites, whose tail is the thin one, hit in $0.833$ of
replicates, with mean absolute error $1.92$ batches.
The cost is a false alarm rate of $0.375$ on Income clean sites and
$0.125$ on Public Coverage clean sites: any crossing counts, whether or
not it falls near $t^\star$.
Among the Income clean curves that do cross, the mean absolute error is
$4.17$ batches, so the typical alarm is not a call at the change.
A threshold of $0.01$ above a two-point calibration is enough to rank
sites when the window is the full recent block, and it is not enough to
call a changepoint on a short window without a second screen.
The noise site crosses the threshold in two of four Income replicates and
in one of four Public Coverage replicates.
One Income crossing lands on batch $10$, inside the two-batch window; the
other is six batches later.
The Public Coverage crossing is three batches after $t^\star$ and is not
counted as a hit.

\begin{table}[t]
\centering
\caption{Time localization over four seeds.
Online hit rate is the fraction of sites whose first threshold crossing
lies within two batches of $t^\star$.
False-alarm rate is computed only for clean sites.
Derivative MAE is the absolute error of the largest first difference.}
\label{tab:time}
\begin{tabular}{llcccc}
\toprule
Task & Type & Online hit & False alarm & Online MAE & Derivative MAE \\
\midrule
Income & Clean & 0.062 & 0.375 & 4.17 & 4.06 \\
Income & Linear & 0.833 & --- & 1.92 & 2.83 \\
Income & Nonlinear & 1.000 & --- & 1.62 & 3.00 \\
Income & Noise & 0.250 & --- & 4.00 & 3.25 \\
Public Coverage & Clean & 0.062 & 0.125 & 3.00 & 3.88 \\
Public Coverage & Linear & 1.000 & --- & 1.42 & 3.00 \\
Public Coverage & Nonlinear & 1.000 & --- & 1.50 & 3.00 \\
Public Coverage & Noise & 0.000 & --- & 3.00 & 4.00 \\
\bottomrule
\end{tabular}
\end{table}

\subsection{Concept-group type resolution}
\label{sec:type}

Detecting a primary source does not say whether the new concept is the
linear halfspace or one of the two nonlinear rules.
The dictionary is fixed before looking at site labels.
On a window with outcome $y$, three logistic regressions are cross-fit
against an intercept-only Brier score:
$(Z_1,Z_2)$, the scalar $Z_1 Z_2$, and $(\sin Z_1,\, Z_3^2)$.
Group importance is the drop in Brier.
The call is the maximizing group.
One call uses every row in the recent window.
The other uses only the upper tail of $X_1$, the region in which the
linear rule was planted.
Accuracy is computed on the five structured sites only.
Top-five precision asks a different question: of the five sites with the
largest $\Delta^{\mathrm{cf}}$, how many lie in $\mathcal{P}$.

\begin{table}[t]
\centering
\caption{Type resolution at $\alpha=1$, abrupt activation, and $q=0.5$.
Accuracy is the fraction of structured sites whose maximizing concept group
equals the planted mechanism, averaged over four seeds.
Recall is pooled over both tasks and all four seeds, hence over $24$ linear
sites and $8$ sites of each nonlinear mechanism.
Top-five precision of $\Delta^{\mathrm{cf}}$ is $1$ in every seed on both
tasks.}
\label{tab:type}
\begin{tabular}{llcc}
\toprule
Call & Mechanism or task & Accuracy or recall \\
\midrule
All rows & Income, accuracy & 0.700 \\
All rows & Public Coverage, accuracy & 0.700 \\
Upper tail & Income, accuracy & 0.750 \\
Upper tail & Public Coverage, accuracy & 0.900 \\
\midrule
All rows & Linear, recall & 0.500 \\
All rows & Product, recall & 1.000 \\
All rows & Sine-square, recall & 1.000 \\
Upper tail & Linear, recall & 0.917 \\
Upper tail & Product, recall & 0.375 \\
Upper tail & Sine-square, recall & 1.000 \\
\bottomrule
\end{tabular}
\end{table}

Table~\ref{tab:type} separates detection from typing.
Top-five precision is $1$ in every seed: the five largest values of
$\Delta^{\mathrm{cf}}$ are exactly the five structured sites.
Conditional on that detection, the full-window dictionary recovers both
nonlinear mechanisms in every replicate and recovers the linear halfspace
in half of the linear sites.
The miss is structural.
Outside the tail the linear site still carries $y^{\mathrm{base}}$, and on
Income that baseline outcome is already monotone in hours and schooling.
An additive logistic fit can look competitive with the product and
sine-square groups even when the rewritten slice is a halfspace in
$(Z_1,Z_2)$.
Restricting the call to the upper tail raises linear recall from $0.500$
to $0.917$ and raises Public Coverage accuracy from $0.700$ to $0.900$.
The same restriction lowers product recall from $1$ to $0.375$.
The product rule was planted on every row.
Inside the upper tail the factors are truncated, and the additive group
in $(Z_1,Z_2)$ is no longer dominated by the product scalar.
The sine-square group remains separated on either window.
Type resolution is therefore a property of the region on which the group
is scored, not a single accuracy number.

\subsection{Feature support and one-shot aggregation}
\label{sec:features}

Feature scores are computed on the recent window of each structured site
and compared with that site's planted support by AUC.
Four scores are used.
Shift permutation importance is the drop in the squared concept contrast
when a coordinate is permuted inside five quantile bins of the covariate
it is most correlated with.
Both the local forest and the reference forest are re-evaluated on that
same permuted matrix.
The local forest here is fit on the full recent window; the cross-fit in
\eqref{eq:delta-cf} is not reused.
Predictive permutation importance is the drop in held-out AUC for
predicting the recent label.
The difference of that importance between the recent window and the burn-in
is the fourth score's companion; the fourth score itself is the absolute coefficient vector of an
$\ell_1$-penalized logistic regression on the recent window, with
$C=0.5$.

\begin{table}[t]
\centering
\caption{Mean AUC of a feature score against the planted support.
Linear and product supports are $\{X_1,X_2\}$.
The sine-square support is $\{X_1,X_3\}$.
Averages run over the structured sites of four seeds.}
\label{tab:feature}
\begin{tabular}{llccc}
\toprule
Task & Type & Shift permutation & Predictive permutation & $\ell_1$ coefficient \\
\midrule
Income & Linear & 0.906 & 0.927 & 0.786 \\
Income & Nonlinear & 0.789 & 1.000 & 0.930 \\
Public Coverage & Linear & 0.953 & 0.995 & 0.841 \\
Public Coverage & Nonlinear & 0.985 & 1.000 & 0.879 \\
\bottomrule
\end{tabular}
\end{table}

Table~\ref{tab:feature} says that the planted coordinates are recoverable
from a single drifted site.
Predictive permutation AUC is at least $0.927$.
Shift permutation AUC is at least $0.906$ on linear sites and falls to
$0.789$ on Income nonlinear sites, where the product and the sine-square
rule are functions of coordinates that also predict the baseline income
label.
The $\ell_1$ coefficient is weaker than the forest on linear Public
Coverage ($0.841$) and stronger than shift permutation on Income nonlinear
sites ($0.930$).
None of these AUCs is a site-ranking result.
They are the contents of the vector a site would upload.

Aggregation treats each upload as a point on the simplex.
Negative shift-permutation coordinates are truncated at zero and the vector
is renormalized.
The target vector puts mass on a coordinate proportional to the number of
structured sites that plant it.
Uniform averaging, the coordinate-wise median, inverse-variance weights
with variance estimated from the permutation repeats and floored at
$10^{-4}$, and weights proportional to $\Delta^{\mathrm{cf}}$ are compared
by AUC, precision at the number of positive target coordinates, and
Spearman correlation with the target.
Leaving the noise site out of the uniform average, and averaging only the
five structured sites, are oracles used as ceilings rather than as methods.

On Public Coverage every shift aggregator has AUC at least $0.984$, and
the $\Delta^{\mathrm{cf}}$-weighted average and the average that drops the
noise site both have AUC $1$ and precision $1$.
On Income the uniform shift average has AUC $0.869$ and precision $0.667$.
Weighting by $\Delta^{\mathrm{cf}}$ raises those figures to $0.917$ and
$0.750$.
The predictive permutation upload and the $\ell_1$ upload, averaged
uniformly, reach AUC $0.988$ on Income.
That is not evidence against the concept score.
On Income the planted coordinates are hours, schooling, and age, which are
also the coordinates that predict $y^{\mathrm{base}}$ on a clean site.
A clean site's predictive vector already points at the target.
The shift vector is the one that is trying to upload the change rather than
the baseline regression, and it is the vector that improves when clean and
noise sites are down-weighted by $\Delta^{\mathrm{cf}}$.

Leave-one-site-out influence, defined as one minus the Spearman
correlation between the uniform aggregate and the aggregate with one site
removed, does not identify primary sources.
Mean influence on Income is $0.033$ for clean sites, $0.013$ for linear
sites, $0.018$ for the noise site, and $0.021$ for nonlinear sites.
On Public Coverage every type mean lies between $0.005$ and $0.008$.
The uploaded vectors agree with each other closely enough that deleting one
site barely moves the ranking.
Site-level $\Delta^{\mathrm{cf}}$ is the functional that separates primary
sources from clean sites.
Influence on a simplex average is not.

\subsection{Intensity, tail mass, and gradual activation}
\label{sec:sensitivity}

The main grid holds $\alpha=1$, abrupt activation, and $q=0.5$.
The sensitivity grid varies those three levers on Income only, with three
seeds and the same score, and it does not recompute feature uploads or the
joint two-sample test.
Because the $0.5$ and $0.7$ quantiles of hours worked are the same mass
point at $40$ hours, $q=0.7$ selects the same tail as $q=0.5$ and is
omitted.
The quantile $0.8$ moves the threshold to $45$ hours and leaves about
$17\%$ of rows in the linear tail.
The quantile $0.9$ moves it to $50$ hours and leaves about $8\%$.
Nonlinear sites continue to rewrite every activated row.
Only the linear sites become thinner.

\begin{table}[t]
\centering
\caption{Income concordance $C(\Delta^{\mathrm{cf}};\mathcal{P},\mathcal{C})$
under weaker rewrites and thinner linear tails.
Three seeds.
The worst seed is in parentheses.
Quantile $0.5$ is the main tail, about $23\%$ of rows because of the mass
at $40$ hours.}
\label{tab:sensitivity}
\begin{tabular}{lcccccc}
\toprule
& \multicolumn{2}{c}{$q=0.5$} & \multicolumn{2}{c}{$q=0.8$} & \multicolumn{2}{c}{$q=0.9$} \\
\cmidrule(lr){2-3}\cmidrule(lr){4-5}\cmidrule(lr){6-7}
$\alpha$ & Abrupt & Gradual & Abrupt & Gradual & Abrupt & Gradual \\
\midrule
1.0 & 1.000 {\scriptsize(1.000)} & 1.000 {\scriptsize(1.000)} & 1.000 {\scriptsize(1.000)} & 1.000 {\scriptsize(1.000)} & 0.850 {\scriptsize(0.750)} & 0.867 {\scriptsize(0.800)} \\
0.7 & 1.000 {\scriptsize(1.000)} & 1.000 {\scriptsize(1.000)} & 0.967 {\scriptsize(0.950)} & 0.950 {\scriptsize(0.900)} & 0.817 {\scriptsize(0.750)} & 0.683 {\scriptsize(0.500)} \\
0.4 & 0.783 {\scriptsize(0.650)} & 0.750 {\scriptsize(0.650)} & 0.617 {\scriptsize(0.450)} & 0.600 {\scriptsize(0.450)} & 0.550 {\scriptsize(0.450)} & 0.550 {\scriptsize(0.450)} \\
\bottomrule
\end{tabular}
\end{table}

\begin{figure}[t]
\centering
\includegraphics[width=\textwidth]{sensitivity_concordance.png}
\caption{Income concordance of $\Delta^{\mathrm{cf}}$ against clean sites.
Each panel is one linear-tail quantile.
Cell entries are means over three seeds.
Means stay at or above $0.8$ for $\alpha\in\{0.7,1\}$ down to a linear tail
of about $17\%$ of rows.
They fall below $0.8$ when $\alpha=0.4$, and when a tail of about $8\%$ is
combined with $\alpha=0.7$ and a gradual ramp.
Full intensity on that thin tail stays above $0.8$ on average and not in
every seed; the worst seeds are in Table~\ref{tab:sensitivity}.}
\label{fig:sensitivity}
\end{figure}

Table~\ref{tab:sensitivity} and Figure~\ref{fig:sensitivity} locate the
boundary of the $0.8$ target.
At full intensity the concordance remains $1$ down to a linear tail of
about $17\%$ of rows, whether the change is abrupt or spread over four
batches.
At a tail of about $8\%$, full abrupt intensity falls to mean concordance
$0.850$ with a worst seed of $0.750$, so the mean still clears $0.8$ and
the guarantee across seeds does not.
Gradual activation at that thin tail has the same character
($0.867$, worst seed $0.800$).
Lowering the rewrite probability to $0.7$ preserves the target on the main
tail and on the $17\%$ tail.
Combining $\alpha=0.7$ with the $8\%$ tail and a four-batch ramp drops the
concordance to $0.683$, with a worst seed of $0.500$.
At $\alpha=0.4$ every column is below $0.8$.
The main-tail abrupt entry is $0.783$ with a worst seed of $0.650$.
On that same $\alpha=0.4$ slice, the whole-window score without the tail
scan has concordance $0.433$.
The region maximum is what keeps a weak but structured rewrite near the
target.
It does not keep an arbitrarily weak rewrite above the target.

\subsection{What the benchmark supports}
\label{sec:support}

The benchmark supports four claims, each tied to a table.
First, when the rewrite is complete and the linear tail contains about a
fifth of the rows or more, $\Delta^{\mathrm{cf}}$ orders every structured
site ahead of every clean site and ahead of the noise site, on two real
covariate tables, in every seed
(Table~\ref{tab:concordance}).
Second, the functionals that ignore either the existence of a local concept
or the region in which that concept lives do not.
Label total variation and the covariate arms are near chance against clean
sites.
The joint two-sample test and the pooled Brier excess promote the noise
site, completely so for linear sites under the pooled excess on Income.
Third, the time profile of the same score locates an abrupt change to
within about two batches once the scoring window is taken into account,
with a false-alarm rate on clean Income sites that a deployment threshold
would still have to budget for
(Table~\ref{tab:time}).
Fourth, a pre-specified concept-group dictionary can name the nonlinear
mechanisms from the full window and the linear mechanism from the tail, and
it cannot do both from a single region
(Table~\ref{tab:type}).

The benchmark does not support a claim that leave-one-site-out influence
on an averaged variable-importance vector identifies the primary site.
Those vectors are too similar across sites for the deletion of one upload
to move the rank correlation.
It also does not support a claim at very small intensity.
Below a rewrite probability of one half, or inside a linear tail that
covers less than a tenth of the window, pairwise concordance falls through
$0.8$ on at least one seed.
Both limits are visible in the same design that produces the perfect
separation at full intensity, which is the reason for reporting them beside
the headline rather than as a separate simulation.

% after defining nothing: the section is self-contained.

\section{Primary-source identification across heterogeneous sites}
\label{sec:primary-source}

The manuscript separates covariate shift, a change in $P(X)$, from concept
drift, a change in $P(Y\mid X)$ between a reference batch
$D_{\mathrm{exist}}$ and a new batch $D_{\mathrm{new}}$.
Feature-level procedures in the paper, including the pseudo-outcome
risk and conditional permutation importance, score coordinates inside one
batch.
The federated aggregation in Appendix~11 then asks whether those
variable-importance vectors can be merged without moving raw rows.
That merge treats every site as a noisy look at one common shift.
A prior operational question is which sites are the primary sources of the
shift, and which sites are either stable or actively destroying the
consensus.
This section constructs a ten-site benchmark in which the answer is known,
and evaluates a counterfactual concept score against that answer.

The design isolates concept drift.
Every site is a random slice of one real covariate table, so the marginal
distribution $P(X)$ is shared.
The American Community Survey extracts below are cross-sections, not
panels, so the batch index is a seeded random order rather than a clock.
Stationarity of $P(X)$ is then a construction, and a covariate arm that
fires is evidence of a bug in the generator rather than a feature of the
task.
Labels follow the real outcome through a burn-in and are then rewritten,
on a known subset of sites, by a known rule.

\subsection{Data and the ten-site generator}
\label{sec:dgp}

Two WhyShift tasks supply the base table $(X,y^{\mathrm{base}})$, both
taken from the 2018 ACS one-year person file for California and both using
the official \texttt{folktables} feature lists.
ACS Income has $p=10$ covariates and the indicator that personal income
exceeds $\$50{,}000$.
ACS Public Coverage has $p=19$ covariates and the indicator of public
health-insurance coverage.
In each Monte Carlo replicate the analysis sample is winsorized at the
first and ninety-ninth percentiles, standardized, and cut into $K=10$
non-overlapping sites of $n_k=1600$ rows.
Each site is then cut into $T=16$ batches of $100$ rows.
The burn-in is the first four batches.
The concept change, when it is switched on, begins at batch
$t^\star=8$.

The ten sites have four roles.
Four clean sites keep $y=y^{\mathrm{base}}$ at every batch.
Three linear sites rewrite labels inside an upper tail of a designated
coordinate $X_1$.
Two nonlinear sites rewrite every activated row.
One noise site replaces activated labels with independent
$\mathrm{Bernoulli}(1/2)$ draws.
The noise site is a consensus breaker.
It is not a primary source: there is no new concept for a local model to
learn.
Table~\ref{tab:roles} records which coordinates carry the rules.
Income uses hours worked, educational attainment, and age.
Public Coverage uses age, educational attainment, and personal income.
These are the continuous or ordinal columns in each official feature list.
The remaining columns stay in the table as nuisance covariates.

\begin{table}[t]
\centering
\caption{Rule coordinates. $Z_j$ denotes the analysis-sample standardize
of the named column. Linear sites rewrite only the upper tail of $Z_1$.
Nonlinear sites rewrite every activated row.}
\label{tab:roles}
\begin{tabular}{llccc}
\toprule
Task & Role of $(X_1,X_2,X_3)$ & $X_1$ & $X_2$ & $X_3$ \\
\midrule
ACS Income & hours, schooling, age & WKHP & SCHL & AGEP \\
ACS Public Coverage & age, schooling, income & AGEP & SCHL & PINCP \\
\bottomrule
\end{tabular}
\end{table}

Write $Z$ for the standardized covariates and let
$M_k(q)=\{Z_1>Q_q(Z_1)\}$ be the upper tail at quantile $q$, computed on
the pooled analysis sample.
The main design uses $q=0.5$.
On a linear site the candidate label is the halfspace
\begin{equation}
\label{eq:linear-rule}
y^{\mathrm{rule}}
=
\mathbf{1}\!\left\{
\sigma_k(\beta_{k1} Z_1+\beta_{k2} Z_2)>c_k
\right\}
\quad\text{on }M_k(q),
\end{equation}
and $y^{\mathrm{base}}$ off that tail.
The three coefficient vectors are
$(1,0.35)$, $(0.25,1)$, and $(-0.85,0.55)$.
The threshold $c_k$ is the quantile of the linear score that matches the
positive rate of $y^{\mathrm{base}}$ on $M_k(q)$, and the sign
$\sigma_k\in\{\pm 1\}$ is the orientation whose matched labels agree less
with $y^{\mathrm{base}}$.
The sign choice is part of the generator.
It keeps the planted concept from collapsing back onto the real outcome
when hours or schooling already predict income.
It is not estimated by the detector.

On a nonlinear site the same rate match is applied to every row, not only
the tail:
\begin{equation}
\label{eq:nonlinear-rule}
y^{\mathrm{rule}}
=
\mathbf{1}\{Z_1 Z_2>c\}
\qquad\text{or}\qquad
y^{\mathrm{rule}}
=
\mathbf{1}\{\sin(Z_1)+Z_3^2>c\}.
\end{equation}
The first site uses the product and the second uses the sine-square rule.
The planted feature support is $\{X_1,X_2\}$ for every linear site and for
the product site, and $\{X_1,X_3\}$ for the sine-square site.

Activation is a batch weight $w_t\in[0,1]$.
Under abrupt activation, $w_t=\mathbf{1}\{t\ge t^\star\}$.
Under gradual activation, $w_t=0$ before $t^\star$ and then rises linearly
to one over four batches.
Each post-burn-in row is rewritten with probability $\alpha w_t$, and
otherwise keeps $y^{\mathrm{base}}$.
The main grid uses $\alpha=1$.
The noise site uses the same weights and then draws a fresh coin flip.
Clean sites ignore the weight.

Two consequences are worth stating before any score is computed.
First, because the structured rules match the local positive rate, a
marginal label discrepancy is not a substitute for concept drift.
The largest average gap between the post-change positive rate and the
baseline, among structured sites, is $2.4$ percentage points
(Income, nonlinear).
Second, hours worked in the ACS piles up at $40$.
The empirical quantiles of WKHP at $0.5$ and at $0.7$ are both $40$, so
$\{Z_1>Q_{0.5}\}$ and $\{Z_1>Q_{0.7}\}$ are the same set, about $23\%$ of
Income rows.
The median tail in the main Income design is therefore thinner than a
continuous coordinate would suggest.
Public Coverage uses age as $X_1$, and that tail is about half the sample.
Table~\ref{tab:dgp-check} records the realized flip rates.
Structured sites move roughly half the labels inside the region they
rewrite, and they leave the marginal positive rate near the baseline.
The noise site flips about half of all post-burn-in labels and shifts the
positive rate to one half.

\begin{table}[t]
\centering
\caption{Realized label changes after $t^\star$, averaged over four seeds
at $\alpha=1$, abrupt activation, and $q=0.5$.
Flip rate is the fraction of post-burn-in rows whose label differs from
$y^{\mathrm{base}}$.
Agreement is the fraction of the rewritten region on which the rate-matched
rule reproduces $y^{\mathrm{base}}$.
Clean sites are not rewritten, so agreement is one.
The noise site has no rule, so agreement is left blank.}
\label{tab:dgp-check}
\begin{tabular}{llcccc}
\toprule
Task & Site type & Flip rate & Rule agreement & Post positive rate & Base positive rate \\
\midrule
Income & Clean & 0.000 & 1.000 & 0.406 & 0.406 \\
Income & Linear & 0.131 & 0.427 & 0.400 & 0.414 \\
Income & Nonlinear & 0.555 & 0.436 & 0.396 & 0.420 \\
Income & Noise & 0.498 & --- & 0.491 & 0.410 \\
Public Coverage & Clean & 0.000 & 1.000 & 0.369 & 0.367 \\
Public Coverage & Linear & 0.259 & 0.476 & 0.378 & 0.372 \\
Public Coverage & Nonlinear & 0.492 & 0.506 & 0.371 & 0.374 \\
Public Coverage & Noise & 0.496 & --- & 0.495 & 0.362 \\
\bottomrule
\end{tabular}
\end{table}

\subsection{The counterfactual concept score}
\label{sec:delta-cf}

Let $\mathcal{D}_0$ be the pooled burn-in, eight thousand rows in the main
design, still generated from $y^{\mathrm{base}}$ at every site.
Let $\mathcal{D}_k$ be the most recent six batches at site $k$, six hundred
rows, all of them inside the drift window when activation is abrupt.
A random forest with $25$ trees, depth $8$, and minimum leaf size $5$ is
the outcome model throughout, in line with the recommendation that the
forest be used for distance estimation rather than for a causal reading of
the meta-learner.
Local predictions on $\mathcal{D}_k$ are three-fold cross-fit probabilities
$\mu_k^{\mathrm{oof}}$.
The reference probabilities $\mu_0^{\mathrm{sm}}$ are an average of two
forests, each trained on a random subset of $\mathcal{D}_0$ whose size
equals the local training set, two thirds of the current window.
Matching the training size stops the larger reference sample from looking
like a concept change.
A third forest trained on all of $\mathcal{D}_0$ is retained only as the
pooled excess defined below.
It is not the headline score.

On a region $\mathcal{R}$ that has both classes in the current window and
in the reference, at least fifty current rows, and at least eighty
reference rows, write $B(y,p)$ for the Brier score and
\begin{equation}
\label{eq:delta-region}
\Delta_k(\mathcal{R})
=
\Big(B(\mu_0^{\mathrm{sm}})-B(\mu_k^{\mathrm{oof}})\Big)_+
\cdot
\Big(2\big(\mathrm{AUC}(\mu_k^{\mathrm{oof}})-{\textstyle\frac12}\big)\Big)_+.
\end{equation}
The first factor is the excess risk of the reference concept over a local
concept of the same capacity.
The second factor is zero when the local model does not beat a coin, and
it equals one when the local AUC equals one.
The headline score is the maximum over a dictionary that does not use
labels,
\begin{equation}
\label{eq:delta-cf}
\Delta_k^{\mathrm{cf}}
=
\max_{\mathcal{R}\in\{\mathrm{all},\,\mathrm{high},\,\mathrm{low}\}}
\Delta_k(\mathcal{R}),
\end{equation}
where ``high'' and ``low'' are the two sides of the burn-in median of
$X_1$.
The median is computed on $\mathcal{D}_0$ only.
The whole-window score $\Delta_k(\mathrm{all})$ is reported beside the
maximum so that the cost of ignoring the tail is visible.

The same two probability vectors define the concept contrast
\begin{equation}
\label{eq:tau}
\tau_k^2
=
\frac{1}{|\mathcal{D}_k\cap\mathcal{R}_k^\star|}
\sum_{i\in\mathcal{D}_k\cap\mathcal{R}_k^\star}
\big(\mu_k^{\mathrm{oof}}(X_i)-\mu_0^{\mathrm{sm}}(X_i)\big)^2,
\end{equation}
with $\mathcal{R}_k^\star$ the region attaining the maximum in
\eqref{eq:delta-cf}.
This is the concept contrast that the pseudo-outcome decomposition isolates
from the covariate arm once $P(X)$ is shared across $D_{\mathrm{exist}}$
and $D_{\mathrm{new}}$.
The benchmark does not re-estimate a propensity: under the common
covariate table that weight is not where the signal sits.
Gating by the same skill factor as in \eqref{eq:delta-region} produces
$\tau_k^{2,\mathrm{g}}$.
Raw $\tau_k^2$ is retained because a flat local predictor on a noise site
can still sit far from a sharp reference predictor.
The gate is the part of $\Delta_k^{\mathrm{cf}}$ that refuses to call that
gap a primary source.

For a score $s$ and two disjoint collections of sites $\mathcal{A}$ and
$\mathcal{B}$, the pairwise concordance is
\begin{equation}
\label{eq:concordance}
C(s;\mathcal{A},\mathcal{B})
=
\frac{1}{|\mathcal{A}|\,|\mathcal{B}|}
\sum_{a\in\mathcal{A}}\sum_{b\in\mathcal{B}}
\mathbf{1}\{s_a>s_b\}.
\end{equation}
Ties count as failures.
The primary set $\mathcal{P}$ is the five structured sites.
The target for the headline score is
$C(\Delta^{\mathrm{cf}};\mathcal{P},\mathcal{C})>0.8$,
where $\mathcal{C}$ is the set of four clean sites.
The same functional is reported for $\mathcal{P}$ against the noise site,
and separately for the linear and nonlinear subsets of $\mathcal{P}$.
Every entry below is a mean over four seeds, with the worst seed in
parentheses, unless a row is explicitly a per-type mean.

Negative controls use the same windows and never see the label rule.
The domain AUC is the cross-fit AUC of a classifier trained to separate
burn-in rows from recent rows using $X$ alone, once with the forest and
once with logistic regression.
The covariate discrepancy is an RBF MMD on subsamples of size $250$, with
bandwidth equal to the median of the squared Euclidean distances.
A joint classifier two-sample test appends $y$ to $X$ before the same
forest AUC is computed.
Marginal label total variation is the absolute difference of positive rates
between the recent window and the burn-in.

\begin{table}[t]
\centering
\caption{Mean score by site type at $\alpha=1$, abrupt activation, and
$q=0.5$. Four seeds. $\Delta^{\mathrm{cf}}$ is \eqref{eq:delta-cf}.
$\Delta(\mathrm{all})$ is the whole-window score.
$\Delta(\mathrm{high})$ is the upper tail of $X_1$.
Domain AUC is the logistic classifier on $X$ only.}
\label{tab:means}
\small
\begin{tabular}{llrrrrrr}
\toprule
Task & Type & $\Delta^{\mathrm{cf}}$ & $\Delta(\mathrm{all})$ & $\Delta(\mathrm{high})$ & $\tau^2$ & Skill & Domain AUC \\
\midrule
Income & Clean & 0.003 & 0.000 & 0.002 & 0.015 & 0.696 & 0.498 \\
Income & Linear & 0.253 & 0.024 & 0.253 & 0.181 & 0.937 & 0.482 \\
Income & Nonlinear & 0.336 & 0.298 & 0.316 & 0.313 & 0.897 & 0.494 \\
Income & Noise & 0.001 & 0.000 & 0.000 & 0.096 & 0.008 & 0.494 \\
Public Coverage & Clean & 0.001 & 0.000 & 0.001 & 0.014 & 0.364 & 0.503 \\
Public Coverage & Linear & 0.200 & 0.052 & 0.200 & 0.132 & 0.943 & 0.493 \\
Public Coverage & Nonlinear & 0.236 & 0.182 & 0.179 & 0.159 & 0.970 & 0.478 \\
Public Coverage & Noise & 0.003 & 0.002 & 0.003 & 0.046 & 0.089 & 0.528 \\
\bottomrule
\end{tabular}
\end{table}

Table~\ref{tab:means} is the raw material of the concordance claim.
On both tasks, and for every linear site, the maximum is exactly the
upper-tail score: $\Delta^{\mathrm{cf}}=\Delta(\mathrm{high})$.
On Income that tail score is $0.253$ against a whole-window score of
$0.024$.
The whole window is mostly unrewritten rows, and a single Brier average
dilutes a real halfspace by about an order of magnitude.
The noise skill is $0.008$ on Income and $0.089$ on Public Coverage, so
the gate shuts even though the ungated contrast $\tau^2$ on Income noise
($0.096$) exceeds the clean contrast ($0.015$).
Logistic domain AUCs sit on $1/2$ for every type.
The largest is $0.528$, on the Public Coverage noise site.
The forest domain AUC, not shown column by column, has type means between
$0.490$ and $0.514$, and equals $0.514$ on that same noise site.
Neither is evidence of a covariate shift.

\subsection{Concordance}
\label{sec:concordance}

\begin{table}[t]
\centering
\caption{Pairwise concordance \eqref{eq:concordance} at $\alpha=1$, abrupt
activation, and $q=0.5$.
Each cell is the mean over four seeds, with the worst seed in parentheses.
$\mathcal{P}$ is the five structured sites and $\mathcal{C}$ is the four
clean sites.}
\label{tab:concordance}
\small
\begin{tabular}{llcccc}
\toprule
Task & Score & $\mathcal{P}>\mathcal{C}$ & $\mathcal{P}>\mathrm{noise}$ & Linear $>\mathcal{C}$ & Nonlinear $>\mathcal{C}$ \\
\midrule
Income & $\Delta^{\mathrm{cf}}$ & 1.000 {\scriptsize(1.000)} & 1.000 {\scriptsize(1.000)} & 1.000 {\scriptsize(1.000)} & 1.000 {\scriptsize(1.000)} \\
Income & $\Delta(\mathrm{all})$ & 1.000 {\scriptsize(1.000)} & 1.000 {\scriptsize(1.000)} & 1.000 {\scriptsize(1.000)} & 1.000 {\scriptsize(1.000)} \\
Income & $\tau^{2,\mathrm{g}}$ & 1.000 {\scriptsize(1.000)} & 1.000 {\scriptsize(1.000)} & 1.000 {\scriptsize(1.000)} & 1.000 {\scriptsize(1.000)} \\
Income & Pareto rank & 1.000 {\scriptsize(1.000)} & 1.000 {\scriptsize(1.000)} & 1.000 {\scriptsize(1.000)} & 1.000 {\scriptsize(1.000)} \\
Income & $\tau^2$ & 1.000 {\scriptsize(1.000)} & 1.000 {\scriptsize(1.000)} & 1.000 {\scriptsize(1.000)} & 1.000 {\scriptsize(1.000)} \\
Income & Pooled excess & 1.000 {\scriptsize(1.000)} & 0.400 {\scriptsize(0.400)} & 1.000 {\scriptsize(1.000)} & 1.000 {\scriptsize(1.000)} \\
Income & Label TV & 0.500 {\scriptsize(0.100)} & 0.050 {\scriptsize(0.000)} & 0.500 {\scriptsize(0.083)} & 0.500 {\scriptsize(0.125)} \\
Income & Joint C2ST & 0.850 {\scriptsize(0.800)} & 0.350 {\scriptsize(0.200)} & 0.750 {\scriptsize(0.667)} & 1.000 {\scriptsize(1.000)} \\
Income & RF domain AUC & 0.487 {\scriptsize(0.400)} & 0.500 {\scriptsize(0.000)} & 0.521 {\scriptsize(0.333)} & 0.438 {\scriptsize(0.250)} \\
Income & RBF MMD & 0.412 {\scriptsize(0.300)} & 0.500 {\scriptsize(0.000)} & 0.375 {\scriptsize(0.167)} & 0.469 {\scriptsize(0.125)} \\
\midrule
Public Coverage & $\Delta^{\mathrm{cf}}$ & 1.000 {\scriptsize(1.000)} & 1.000 {\scriptsize(1.000)} & 1.000 {\scriptsize(1.000)} & 1.000 {\scriptsize(1.000)} \\
Public Coverage & $\Delta(\mathrm{all})$ & 1.000 {\scriptsize(1.000)} & 1.000 {\scriptsize(1.000)} & 1.000 {\scriptsize(1.000)} & 1.000 {\scriptsize(1.000)} \\
Public Coverage & $\tau^{2,\mathrm{g}}$ & 1.000 {\scriptsize(1.000)} & 1.000 {\scriptsize(1.000)} & 1.000 {\scriptsize(1.000)} & 1.000 {\scriptsize(1.000)} \\
Public Coverage & Pareto rank & 1.000 {\scriptsize(1.000)} & 1.000 {\scriptsize(1.000)} & 1.000 {\scriptsize(1.000)} & 1.000 {\scriptsize(1.000)} \\
Public Coverage & $\tau^2$ & 1.000 {\scriptsize(1.000)} & 1.000 {\scriptsize(1.000)} & 1.000 {\scriptsize(1.000)} & 1.000 {\scriptsize(1.000)} \\
Public Coverage & Pooled excess & 1.000 {\scriptsize(1.000)} & 1.000 {\scriptsize(1.000)} & 1.000 {\scriptsize(1.000)} & 1.000 {\scriptsize(1.000)} \\
Public Coverage & Label TV & 0.613 {\scriptsize(0.350)} & 0.000 {\scriptsize(0.000)} & 0.604 {\scriptsize(0.417)} & 0.625 {\scriptsize(0.250)} \\
Public Coverage & Joint C2ST & 0.950 {\scriptsize(0.850)} & 0.450 {\scriptsize(0.200)} & 0.917 {\scriptsize(0.750)} & 1.000 {\scriptsize(1.000)} \\
Public Coverage & RF domain AUC & 0.438 {\scriptsize(0.350)} & 0.200 {\scriptsize(0.000)} & 0.396 {\scriptsize(0.167)} & 0.500 {\scriptsize(0.250)} \\
Public Coverage & RBF MMD & 0.600 {\scriptsize(0.500)} & 0.200 {\scriptsize(0.000)} & 0.562 {\scriptsize(0.417)} & 0.656 {\scriptsize(0.375)} \\
\bottomrule
\end{tabular}
\end{table}

Table~\ref{tab:concordance} is the primary result.
On both tasks and in every seed,
$C(\Delta^{\mathrm{cf}};\mathcal{P},\mathcal{C})=1$,
and the same equality holds for linear sites alone, for nonlinear sites
alone, and for $\mathcal{P}$ against the noise site.
The whole-window score and the skill-gated contrast $\tau^{2,\mathrm{g}}$
match this ranking at full intensity.
Figure~\ref{fig:concordance} plots the clean-site comparison.
The covariate rows do not.
Random-forest domain AUC and RBF MMD have concordances between $0.38$ and
$0.66$, with worst seeds as low as $0.17$.
Label total variation is consistent with a matched positive rate: its
concordance against clean sites is $0.50$ on Income and $0.61$ on Public
Coverage, and its concordance against the noise site is $0.05$ and $0.00$.
The noise site is the unique marginal shift, and a marginal functional
ranks it first.

The joint classifier two-sample test sees any change in $(X,Y)$.
It separates $\mathcal{P}$ from $\mathcal{C}$ at $0.85$ on Income and $0.95$
on Public Coverage, which clears $0.8$ on average, but it ranks the noise
site as the stronger departure: concordance of $\mathcal{P}$ against noise
is $0.35$ and $0.45$.
On Income the linear-only comparison for the joint test falls to $0.75$.
A test of the joint distribution is not a test for a learnable new concept.

Pooled excess makes the same mistake in a sharper way, and only on the
task where the reference concept is easy.
It is the Brier gap between the forest trained on the entire burn-in and
the local cross-fit model, before size matching and before the skill gate.
On Income its concordance against the noise site is $0.400$ in every seed.
That fraction is exact: both nonlinear sites beat the noise site, and none
of the three linear sites do.
The pooled Income forest is a sharp predictor of $y^{\mathrm{base}}$.
Evaluated on coin-flip labels its Brier rises by $0.084$ on average, while
the same forest evaluated on a linear tail that covers $23\%$ of the window
rises by only $0.033$.
Size matching does not by itself fix the comparison.
On the region that attains $\Delta^{\mathrm{cf}}$, the size-matched excess
on Income noise remains $0.071$, against $0.002$ on clean sites.
Multiplying by the skill factor is what sends the noise score from $0.071$
to $0.001$.
Public Coverage does not reproduce the $0.40$ figure.
There the baseline concept is harder, the pooled excess on noise is
$0.035$, and every structured site still outranks it.
The failure mode is real on the task where the reference fit is strong,
which is the task on which an ungated risk gap is most tempting.

\begin{figure}[t]
\centering
\includegraphics[width=\textwidth]{concordance_cd_vs_clean.png}
\caption{Pairwise concordance against clean sites.
The dashed line is the $0.8$ target.
$\Delta^{\mathrm{cf}}$, the whole-window score, skill-gated $\tau^2$, raw
$\tau^2$, and the Pareto rank clear the line on both tasks.
The joint classifier two-sample test also clears it on average; its failure
is the noise contrast in Table~\ref{tab:concordance}.
Label total variation, the random-forest domain AUC, and RBF MMD do not
clear the line.
Source: abrupt activation, $\alpha=1$, $q=0.5$, four seeds.}
\label{fig:concordance}
\end{figure}

The Pareto comparison uses three axes when the time curve below is
available, and two otherwise.
The axes are $\Delta^{\mathrm{cf}}$, $\tau^{2,\mathrm{g}}$, and the time
jump, defined as the maximum of the curve on endpoints at or after
$t^\star$ minus the maximum on endpoints before $t^\star$.
A site's rank is one plus the number of sites that are at least as large
on every axis and strictly larger on at least one.
Concordance is computed on the negative rank, so a smaller rank is a
higher score.
Mean ranks on Income are $1.875$ for nonlinear sites, $3.250$ for linear
sites, $6.625$ for clean sites, and $6.750$ for the noise site.
On Public Coverage they are $1.375$, $2.667$, $7.188$, and $6.500$.
The undominated front is mostly the nonlinear sites: half of them on
Income and five eighths on Public Coverage.
No clean site and no noise site is undominated.
Linear sites are often dominated by a nonlinear site, because the
nonlinear rewrite covers every row and the linear rewrite covers a tail,
but they remain strictly ahead of every clean site.
That is the sense in which the Pareto order puts primary sources first.
It does not claim that every primary source lies on the front.

\subsection{Time localization}
\label{sec:time}

The curve at site $k$ reuses the burn-in reference and recomputes
\eqref{eq:delta-cf} on a sliding window of four batches, using only the
whole window and the upper tail.
The horizontal coordinate in Figure~\ref{fig:time} is the last batch
included in the window.
Drift begins at batch $8$.
A window of length four is entirely inside the drift only when its last
batch is at least $11$.
The figure follows that arithmetic.
Means for clean sites and for the noise site stay below $0.01$ on both
sides of $t^\star$.
Linear and nonlinear means leave zero between batches $9$ and $10$ and
level off at batch $11$.
On Income the nonlinear plateau is near $0.32$ and the linear plateau is
near $0.22$.
On Public Coverage the two plateaus sit together near $0.18$, which is
consistent with the similar headline scores in Table~\ref{tab:means}.

\begin{figure}[t]
\centering
\includegraphics[width=\textwidth]{time_profile_delta_cf.png}
\caption{Mean $\Delta^{\mathrm{cf}}$ against the last batch in a window of
length four.
The vertical line is $t^\star=8$.
The curve levels off at batch $11$, the first endpoint whose window lies
entirely after the change.
Clean and noise traces remain near zero.
Source: abrupt activation, $\alpha=1$, $q=0.5$, four seeds.}
\label{fig:time}
\end{figure}

Two detectors are applied to each curve without being told $t^\star$.
The online rule calibrates on the first two evaluated endpoints, sets the
threshold at that maximum plus $0.01$, and reports the first later endpoint
that exceeds the threshold.
A hit is an estimate within two batches of $t^\star$.
The derivative rule reports the endpoint of the largest first difference.
Table~\ref{tab:time} shows that the derivative is estimating the onset of
the plateau rather than the onset of the drift.
Its mean absolute error is about three batches on every structured type,
which is one window length.
The online rule is the onset estimator.
It hits every nonlinear site on both tasks and every linear site on Public
Coverage.
Income linear sites, whose tail is the thin one, hit in $0.833$ of
replicates, with mean absolute error $1.92$ batches.
The cost is a false alarm rate of $0.375$ on Income clean sites and
$0.125$ on Public Coverage clean sites: any crossing counts, whether or
not it falls near $t^\star$.
Among the Income clean curves that do cross, the mean absolute error is
$4.17$ batches, so the typical alarm is not a call at the change.
A threshold of $0.01$ above a two-point calibration is enough to rank
sites when the window is the full recent block, and it is not enough to
call a changepoint on a short window without a second screen.
The noise site crosses the threshold in two of four Income replicates and
in one of four Public Coverage replicates.
One Income crossing lands on batch $10$, inside the two-batch window; the
other is six batches later.
The Public Coverage crossing is three batches after $t^\star$ and is not
counted as a hit.

\begin{table}[t]
\centering
\caption{Time localization over four seeds.
Online hit rate is the fraction of sites whose first threshold crossing
lies within two batches of $t^\star$.
False-alarm rate is computed only for clean sites.
Derivative MAE is the absolute error of the largest first difference.}
\label{tab:time}
\begin{tabular}{llcccc}
\toprule
Task & Type & Online hit & False alarm & Online MAE & Derivative MAE \\
\midrule
Income & Clean & 0.062 & 0.375 & 4.17 & 4.06 \\
Income & Linear & 0.833 & --- & 1.92 & 2.83 \\
Income & Nonlinear & 1.000 & --- & 1.62 & 3.00 \\
Income & Noise & 0.250 & --- & 4.00 & 3.25 \\
Public Coverage & Clean & 0.062 & 0.125 & 3.00 & 3.88 \\
Public Coverage & Linear & 1.000 & --- & 1.42 & 3.00 \\
Public Coverage & Nonlinear & 1.000 & --- & 1.50 & 3.00 \\
Public Coverage & Noise & 0.000 & --- & 3.00 & 4.00 \\
\bottomrule
\end{tabular}
\end{table}

\subsection{Concept-group type resolution}
\label{sec:type}

Detecting a primary source does not say whether the new concept is the
linear halfspace or one of the two nonlinear rules.
The dictionary is fixed before looking at site labels.
On a window with outcome $y$, three logistic regressions are cross-fit
against an intercept-only Brier score:
$(Z_1,Z_2)$, the scalar $Z_1 Z_2$, and $(\sin Z_1,\, Z_3^2)$.
Group importance is the drop in Brier.
The call is the maximizing group.
One call uses every row in the recent window.
The other uses only the upper tail of $X_1$, the region in which the
linear rule was planted.
Accuracy is computed on the five structured sites only.
Top-five precision asks a different question: of the five sites with the
largest $\Delta^{\mathrm{cf}}$, how many lie in $\mathcal{P}$.

\begin{table}[t]
\centering
\caption{Type resolution at $\alpha=1$, abrupt activation, and $q=0.5$.
Accuracy is the fraction of structured sites whose maximizing concept group
equals the planted mechanism, averaged over four seeds.
Recall is pooled over both tasks and all four seeds, hence over $24$ linear
sites and $8$ sites of each nonlinear mechanism.
Top-five precision of $\Delta^{\mathrm{cf}}$ is $1$ in every seed on both
tasks.}
\label{tab:type}
\begin{tabular}{llcc}
\toprule
Call & Mechanism or task & Accuracy or recall \\
\midrule
All rows & Income, accuracy & 0.700 \\
All rows & Public Coverage, accuracy & 0.700 \\
Upper tail & Income, accuracy & 0.750 \\
Upper tail & Public Coverage, accuracy & 0.900 \\
\midrule
All rows & Linear, recall & 0.500 \\
All rows & Product, recall & 1.000 \\
All rows & Sine-square, recall & 1.000 \\
Upper tail & Linear, recall & 0.917 \\
Upper tail & Product, recall & 0.375 \\
Upper tail & Sine-square, recall & 1.000 \\
\bottomrule
\end{tabular}
\end{table}

Table~\ref{tab:type} separates detection from typing.
Top-five precision is $1$ in every seed: the five largest values of
$\Delta^{\mathrm{cf}}$ are exactly the five structured sites.
Conditional on that detection, the full-window dictionary recovers both
nonlinear mechanisms in every replicate and recovers the linear halfspace
in half of the linear sites.
The miss is structural.
Outside the tail the linear site still carries $y^{\mathrm{base}}$, and on
Income that baseline outcome is already monotone in hours and schooling.
An additive logistic fit can look competitive with the product and
sine-square groups even when the rewritten slice is a halfspace in
$(Z_1,Z_2)$.
Restricting the call to the upper tail raises linear recall from $0.500$
to $0.917$ and raises Public Coverage accuracy from $0.700$ to $0.900$.
The same restriction lowers product recall from $1$ to $0.375$.
The product rule was planted on every row.
Inside the upper tail the factors are truncated, and the additive group
in $(Z_1,Z_2)$ is no longer dominated by the product scalar.
The sine-square group remains separated on either window.
Type resolution is therefore a property of the region on which the group
is scored, not a single accuracy number.

\subsection{Feature support and one-shot aggregation}
\label{sec:features}

Feature scores are computed on the recent window of each structured site
and compared with that site's planted support by AUC.
Four scores are used.
Shift permutation importance is the drop in the squared concept contrast
when a coordinate is permuted inside five quantile bins of the covariate
it is most correlated with.
Both the local forest and the reference forest are re-evaluated on that
same permuted matrix.
The local forest here is fit on the full recent window; the cross-fit in
\eqref{eq:delta-cf} is not reused.
Predictive permutation importance is the drop in held-out AUC for
predicting the recent label.
The difference of that importance between the recent window and the burn-in
is the fourth score's companion; the fourth score itself is the absolute coefficient vector of an
$\ell_1$-penalized logistic regression on the recent window, with
$C=0.5$.

\begin{table}[t]
\centering
\caption{Mean AUC of a feature score against the planted support.
Linear and product supports are $\{X_1,X_2\}$.
The sine-square support is $\{X_1,X_3\}$.
Averages run over the structured sites of four seeds.}
\label{tab:feature}
\begin{tabular}{llccc}
\toprule
Task & Type & Shift permutation & Predictive permutation & $\ell_1$ coefficient \\
\midrule
Income & Linear & 0.906 & 0.927 & 0.786 \\
Income & Nonlinear & 0.789 & 1.000 & 0.930 \\
Public Coverage & Linear & 0.953 & 0.995 & 0.841 \\
Public Coverage & Nonlinear & 0.985 & 1.000 & 0.879 \\
\bottomrule
\end{tabular}
\end{table}

Table~\ref{tab:feature} says that the planted coordinates are recoverable
from a single drifted site.
Predictive permutation AUC is at least $0.927$.
Shift permutation AUC is at least $0.906$ on linear sites and falls to
$0.789$ on Income nonlinear sites, where the product and the sine-square
rule are functions of coordinates that also predict the baseline income
label.
The $\ell_1$ coefficient is weaker than the forest on linear Public
Coverage ($0.841$) and stronger than shift permutation on Income nonlinear
sites ($0.930$).
None of these AUCs is a site-ranking result.
They are the contents of the vector a site would upload.

Aggregation treats each upload as a point on the simplex.
Negative shift-permutation coordinates are truncated at zero and the vector
is renormalized.
The target vector puts mass on a coordinate proportional to the number of
structured sites that plant it.
Uniform averaging, the coordinate-wise median, inverse-variance weights
with variance estimated from the permutation repeats and floored at
$10^{-4}$, and weights proportional to $\Delta^{\mathrm{cf}}$ are compared
by AUC, precision at the number of positive target coordinates, and
Spearman correlation with the target.
Leaving the noise site out of the uniform average, and averaging only the
five structured sites, are oracles used as ceilings rather than as methods.

On Public Coverage every shift aggregator has AUC at least $0.984$, and
the $\Delta^{\mathrm{cf}}$-weighted average and the average that drops the
noise site both have AUC $1$ and precision $1$.
On Income the uniform shift average has AUC $0.869$ and precision $0.667$.
Weighting by $\Delta^{\mathrm{cf}}$ raises those figures to $0.917$ and
$0.750$.
The predictive permutation upload and the $\ell_1$ upload, averaged
uniformly, reach AUC $0.988$ on Income.
That is not evidence against the concept score.
On Income the planted coordinates are hours, schooling, and age, which are
also the coordinates that predict $y^{\mathrm{base}}$ on a clean site.
A clean site's predictive vector already points at the target.
The shift vector is the one that is trying to upload the change rather than
the baseline regression, and it is the vector that improves when clean and
noise sites are down-weighted by $\Delta^{\mathrm{cf}}$.

Leave-one-site-out influence, defined as one minus the Spearman
correlation between the uniform aggregate and the aggregate with one site
removed, does not identify primary sources.
Mean influence on Income is $0.033$ for clean sites, $0.013$ for linear
sites, $0.018$ for the noise site, and $0.021$ for nonlinear sites.
On Public Coverage every type mean lies between $0.005$ and $0.008$.
The uploaded vectors agree with each other closely enough that deleting one
site barely moves the ranking.
Site-level $\Delta^{\mathrm{cf}}$ is the functional that separates primary
sources from clean sites.
Influence on a simplex average is not.

\subsection{Intensity, tail mass, and gradual activation}
\label{sec:sensitivity}

The main grid holds $\alpha=1$, abrupt activation, and $q=0.5$.
The sensitivity grid varies those three levers on Income only, with three
seeds and the same score, and it does not recompute feature uploads or the
joint two-sample test.
Because the $0.5$ and $0.7$ quantiles of hours worked are the same mass
point at $40$ hours, $q=0.7$ selects the same tail as $q=0.5$ and is
omitted.
The quantile $0.8$ moves the threshold to $45$ hours and leaves about
$17\%$ of rows in the linear tail.
The quantile $0.9$ moves it to $50$ hours and leaves about $8\%$.
Nonlinear sites continue to rewrite every activated row.
Only the linear sites become thinner.

\begin{table}[t]
\centering
\caption{Income concordance $C(\Delta^{\mathrm{cf}};\mathcal{P},\mathcal{C})$
under weaker rewrites and thinner linear tails.
Three seeds.
The worst seed is in parentheses.
Quantile $0.5$ is the main tail, about $23\%$ of rows because of the mass
at $40$ hours.}
\label{tab:sensitivity}
\begin{tabular}{lcccccc}
\toprule
& \multicolumn{2}{c}{$q=0.5$} & \multicolumn{2}{c}{$q=0.8$} & \multicolumn{2}{c}{$q=0.9$} \\
\cmidrule(lr){2-3}\cmidrule(lr){4-5}\cmidrule(lr){6-7}
$\alpha$ & Abrupt & Gradual & Abrupt & Gradual & Abrupt & Gradual \\
\midrule
1.0 & 1.000 {\scriptsize(1.000)} & 1.000 {\scriptsize(1.000)} & 1.000 {\scriptsize(1.000)} & 1.000 {\scriptsize(1.000)} & 0.850 {\scriptsize(0.750)} & 0.867 {\scriptsize(0.800)} \\
0.7 & 1.000 {\scriptsize(1.000)} & 1.000 {\scriptsize(1.000)} & 0.967 {\scriptsize(0.950)} & 0.950 {\scriptsize(0.900)} & 0.817 {\scriptsize(0.750)} & 0.683 {\scriptsize(0.500)} \\
0.4 & 0.783 {\scriptsize(0.650)} & 0.750 {\scriptsize(0.650)} & 0.617 {\scriptsize(0.450)} & 0.600 {\scriptsize(0.450)} & 0.550 {\scriptsize(0.450)} & 0.550 {\scriptsize(0.450)} \\
\bottomrule
\end{tabular}
\end{table}

\begin{figure}[t]
\centering
\includegraphics[width=\textwidth]{sensitivity_concordance.png}
\caption{Income concordance of $\Delta^{\mathrm{cf}}$ against clean sites.
Each panel is one linear-tail quantile.
Cell entries are means over three seeds.
Means stay at or above $0.8$ for $\alpha\in\{0.7,1\}$ down to a linear tail
of about $17\%$ of rows.
They fall below $0.8$ when $\alpha=0.4$, and when a tail of about $8\%$ is
combined with $\alpha=0.7$ and a gradual ramp.
Full intensity on that thin tail stays above $0.8$ on average and not in
every seed; the worst seeds are in Table~\ref{tab:sensitivity}.}
\label{fig:sensitivity}
\end{figure}

Table~\ref{tab:sensitivity} and Figure~\ref{fig:sensitivity} locate the
boundary of the $0.8$ target.
At full intensity the concordance remains $1$ down to a linear tail of
about $17\%$ of rows, whether the change is abrupt or spread over four
batches.
At a tail of about $8\%$, full abrupt intensity falls to mean concordance
$0.850$ with a worst seed of $0.750$, so the mean still clears $0.8$ and
the guarantee across seeds does not.
Gradual activation at that thin tail has the same character
($0.867$, worst seed $0.800$).
Lowering the rewrite probability to $0.7$ preserves the target on the main
tail and on the $17\%$ tail.
Combining $\alpha=0.7$ with the $8\%$ tail and a four-batch ramp drops the
concordance to $0.683$, with a worst seed of $0.500$.
At $\alpha=0.4$ every column is below $0.8$.
The main-tail abrupt entry is $0.783$ with a worst seed of $0.650$.
On that same $\alpha=0.4$ slice, the whole-window score without the tail
scan has concordance $0.433$.
The region maximum is what keeps a weak but structured rewrite near the
target.
It does not keep an arbitrarily weak rewrite above the target.

\subsection{What the benchmark supports}
\label{sec:support}

The benchmark supports four claims, each tied to a table.
First, when the rewrite is complete and the linear tail contains about a
fifth of the rows or more, $\Delta^{\mathrm{cf}}$ orders every structured
site ahead of every clean site and ahead of the noise site, on two real
covariate tables, in every seed
(Table~\ref{tab:concordance}).
Second, the functionals that ignore either the existence of a local concept
or the region in which that concept lives do not.
Label total variation and the covariate arms are near chance against clean
sites.
The joint two-sample test and the pooled Brier excess promote the noise
site, completely so for linear sites under the pooled excess on Income.
Third, the time profile of the same score locates an abrupt change to
within about two batches once the scoring window is taken into account,
with a false-alarm rate on clean Income sites that a deployment threshold
would still have to budget for
(Table~\ref{tab:time}).
Fourth, a pre-specified concept-group dictionary can name the nonlinear
mechanisms from the full window and the linear mechanism from the tail, and
it cannot do both from a single region
(Table~\ref{tab:type}).

The benchmark does not support a claim that leave-one-site-out influence
on an averaged variable-importance vector identifies the primary site.
Those vectors are too similar across sites for the deletion of one upload
to move the rank correlation.
It also does not support a claim at very small intensity.
Below a rewrite probability of one half, or inside a linear tail that
covers less than a tenth of the window, pairwise concordance falls through
$0.8$ on at least one seed.
Both limits are visible in the same design that produces the perfect
separation at full intensity, which is the reason for reporting them beside
the headline rather than as a separate simulation.

% after defining nothing: the section is self-contained.

\section{Primary-source identification across heterogeneous sites}
\label{sec:primary-source}

The manuscript separates covariate shift, a change in $P(X)$, from concept
drift, a change in $P(Y\mid X)$ between a reference batch
$D_{\mathrm{exist}}$ and a new batch $D_{\mathrm{new}}$.
Feature-level procedures in the paper, including the pseudo-outcome
risk and conditional permutation importance, score coordinates inside one
batch.
The federated aggregation in Appendix~11 then asks whether those
variable-importance vectors can be merged without moving raw rows.
That merge treats every site as a noisy look at one common shift.
A prior operational question is which sites are the primary sources of the
shift, and which sites are either stable or actively destroying the
consensus.
This section constructs a ten-site benchmark in which the answer is known,
and evaluates a counterfactual concept score against that answer.

The design isolates concept drift.
Every site is a random slice of one real covariate table, so the marginal
distribution $P(X)$ is shared.
The American Community Survey extracts below are cross-sections, not
panels, so the batch index is a seeded random order rather than a clock.
Stationarity of $P(X)$ is then a construction, and a covariate arm that
fires is evidence of a bug in the generator rather than a feature of the
task.
Labels follow the real outcome through a burn-in and are then rewritten,
on a known subset of sites, by a known rule.

\subsection{Data and the ten-site generator}
\label{sec:dgp}

Two WhyShift tasks supply the base table $(X,y^{\mathrm{base}})$, both
taken from the 2018 ACS one-year person file for California and both using
the official \texttt{folktables} feature lists.
ACS Income has $p=10$ covariates and the indicator that personal income
exceeds $\$50{,}000$.
ACS Public Coverage has $p=19$ covariates and the indicator of public
health-insurance coverage.
In each Monte Carlo replicate the analysis sample is winsorized at the
first and ninety-ninth percentiles, standardized, and cut into $K=10$
non-overlapping sites of $n_k=1600$ rows.
Each site is then cut into $T=16$ batches of $100$ rows.
The burn-in is the first four batches.
The concept change, when it is switched on, begins at batch
$t^\star=8$.

The ten sites have four roles.
Four clean sites keep $y=y^{\mathrm{base}}$ at every batch.
Three linear sites rewrite labels inside an upper tail of a designated
coordinate $X_1$.
Two nonlinear sites rewrite every activated row.
One noise site replaces activated labels with independent
$\mathrm{Bernoulli}(1/2)$ draws.
The noise site is a consensus breaker.
It is not a primary source: there is no new concept for a local model to
learn.
Table~\ref{tab:roles} records which coordinates carry the rules.
Income uses hours worked, educational attainment, and age.
Public Coverage uses age, educational attainment, and personal income.
These are the continuous or ordinal columns in each official feature list.
The remaining columns stay in the table as nuisance covariates.

\begin{table}[t]
\centering
\caption{Rule coordinates. $Z_j$ denotes the analysis-sample standardize
of the named column. Linear sites rewrite only the upper tail of $Z_1$.
Nonlinear sites rewrite every activated row.}
\label{tab:roles}
\begin{tabular}{llccc}
\toprule
Task & Role of $(X_1,X_2,X_3)$ & $X_1$ & $X_2$ & $X_3$ \\
\midrule
ACS Income & hours, schooling, age & WKHP & SCHL & AGEP \\
ACS Public Coverage & age, schooling, income & AGEP & SCHL & PINCP \\
\bottomrule
\end{tabular}
\end{table}

Write $Z$ for the standardized covariates and let
$M_k(q)=\{Z_1>Q_q(Z_1)\}$ be the upper tail at quantile $q$, computed on
the pooled analysis sample.
The main design uses $q=0.5$.
On a linear site the candidate label is the halfspace
\begin{equation}
\label{eq:linear-rule}
y^{\mathrm{rule}}
=
\mathbf{1}\!\left\{
\sigma_k(\beta_{k1} Z_1+\beta_{k2} Z_2)>c_k
\right\}
\quad\text{on }M_k(q),
\end{equation}
and $y^{\mathrm{base}}$ off that tail.
The three coefficient vectors are
$(1,0.35)$, $(0.25,1)$, and $(-0.85,0.55)$.
The threshold $c_k$ is the quantile of the linear score that matches the
positive rate of $y^{\mathrm{base}}$ on $M_k(q)$, and the sign
$\sigma_k\in\{\pm 1\}$ is the orientation whose matched labels agree less
with $y^{\mathrm{base}}$.
The sign choice is part of the generator.
It keeps the planted concept from collapsing back onto the real outcome
when hours or schooling already predict income.
It is not estimated by the detector.

On a nonlinear site the same rate match is applied to every row, not only
the tail:
\begin{equation}
\label{eq:nonlinear-rule}
y^{\mathrm{rule}}
=
\mathbf{1}\{Z_1 Z_2>c\}
\qquad\text{or}\qquad
y^{\mathrm{rule}}
=
\mathbf{1}\{\sin(Z_1)+Z_3^2>c\}.
\end{equation}
The first site uses the product and the second uses the sine-square rule.
The planted feature support is $\{X_1,X_2\}$ for every linear site and for
the product site, and $\{X_1,X_3\}$ for the sine-square site.

Activation is a batch weight $w_t\in[0,1]$.
Under abrupt activation, $w_t=\mathbf{1}\{t\ge t^\star\}$.
Under gradual activation, $w_t=0$ before $t^\star$ and then rises linearly
to one over four batches.
Each post-burn-in row is rewritten with probability $\alpha w_t$, and
otherwise keeps $y^{\mathrm{base}}$.
The main grid uses $\alpha=1$.
The noise site uses the same weights and then draws a fresh coin flip.
Clean sites ignore the weight.

Two consequences are worth stating before any score is computed.
First, because the structured rules match the local positive rate, a
marginal label discrepancy is not a substitute for concept drift.
The largest average gap between the post-change positive rate and the
baseline, among structured sites, is $2.4$ percentage points
(Income, nonlinear).
Second, hours worked in the ACS piles up at $40$.
The empirical quantiles of WKHP at $0.5$ and at $0.7$ are both $40$, so
$\{Z_1>Q_{0.5}\}$ and $\{Z_1>Q_{0.7}\}$ are the same set, about $23\%$ of
Income rows.
The median tail in the main Income design is therefore thinner than a
continuous coordinate would suggest.
Public Coverage uses age as $X_1$, and that tail is about half the sample.
Table~\ref{tab:dgp-check} records the realized flip rates.
Structured sites move roughly half the labels inside the region they
rewrite, and they leave the marginal positive rate near the baseline.
The noise site flips about half of all post-burn-in labels and shifts the
positive rate to one half.

\begin{table}[t]
\centering
\caption{Realized label changes after $t^\star$, averaged over four seeds
at $\alpha=1$, abrupt activation, and $q=0.5$.
Flip rate is the fraction of post-burn-in rows whose label differs from
$y^{\mathrm{base}}$.
Agreement is the fraction of the rewritten region on which the rate-matched
rule reproduces $y^{\mathrm{base}}$.
Clean sites are not rewritten, so agreement is one.
The noise site has no rule, so agreement is left blank.}
\label{tab:dgp-check}
\begin{tabular}{llcccc}
\toprule
Task & Site type & Flip rate & Rule agreement & Post positive rate & Base positive rate \\
\midrule
Income & Clean & 0.000 & 1.000 & 0.406 & 0.406 \\
Income & Linear & 0.131 & 0.427 & 0.400 & 0.414 \\
Income & Nonlinear & 0.555 & 0.436 & 0.396 & 0.420 \\
Income & Noise & 0.498 & --- & 0.491 & 0.410 \\
Public Coverage & Clean & 0.000 & 1.000 & 0.369 & 0.367 \\
Public Coverage & Linear & 0.259 & 0.476 & 0.378 & 0.372 \\
Public Coverage & Nonlinear & 0.492 & 0.506 & 0.371 & 0.374 \\
Public Coverage & Noise & 0.496 & --- & 0.495 & 0.362 \\
\bottomrule
\end{tabular}
\end{table}

\subsection{The counterfactual concept score}
\label{sec:delta-cf}

Let $\mathcal{D}_0$ be the pooled burn-in, eight thousand rows in the main
design, still generated from $y^{\mathrm{base}}$ at every site.
Let $\mathcal{D}_k$ be the most recent six batches at site $k$, six hundred
rows, all of them inside the drift window when activation is abrupt.
A random forest with $25$ trees, depth $8$, and minimum leaf size $5$ is
the outcome model throughout, in line with the recommendation that the
forest be used for distance estimation rather than for a causal reading of
the meta-learner.
Local predictions on $\mathcal{D}_k$ are three-fold cross-fit probabilities
$\mu_k^{\mathrm{oof}}$.
The reference probabilities $\mu_0^{\mathrm{sm}}$ are an average of two
forests, each trained on a random subset of $\mathcal{D}_0$ whose size
equals the local training set, two thirds of the current window.
Matching the training size stops the larger reference sample from looking
like a concept change.
A third forest trained on all of $\mathcal{D}_0$ is retained only as the
pooled excess defined below.
It is not the headline score.

On a region $\mathcal{R}$ that has both classes in the current window and
in the reference, at least fifty current rows, and at least eighty
reference rows, write $B(y,p)$ for the Brier score and
\begin{equation}
\label{eq:delta-region}
\Delta_k(\mathcal{R})
=
\Big(B(\mu_0^{\mathrm{sm}})-B(\mu_k^{\mathrm{oof}})\Big)_+
\cdot
\Big(2\big(\mathrm{AUC}(\mu_k^{\mathrm{oof}})-{\textstyle\frac12}\big)\Big)_+.
\end{equation}
The first factor is the excess risk of the reference concept over a local
concept of the same capacity.
The second factor is zero when the local model does not beat a coin, and
it equals one when the local AUC equals one.
The headline score is the maximum over a dictionary that does not use
labels,
\begin{equation}
\label{eq:delta-cf}
\Delta_k^{\mathrm{cf}}
=
\max_{\mathcal{R}\in\{\mathrm{all},\,\mathrm{high},\,\mathrm{low}\}}
\Delta_k(\mathcal{R}),
\end{equation}
where ``high'' and ``low'' are the two sides of the burn-in median of
$X_1$.
The median is computed on $\mathcal{D}_0$ only.
The whole-window score $\Delta_k(\mathrm{all})$ is reported beside the
maximum so that the cost of ignoring the tail is visible.

The same two probability vectors define the concept contrast
\begin{equation}
\label{eq:tau}
\tau_k^2
=
\frac{1}{|\mathcal{D}_k\cap\mathcal{R}_k^\star|}
\sum_{i\in\mathcal{D}_k\cap\mathcal{R}_k^\star}
\big(\mu_k^{\mathrm{oof}}(X_i)-\mu_0^{\mathrm{sm}}(X_i)\big)^2,
\end{equation}
with $\mathcal{R}_k^\star$ the region attaining the maximum in
\eqref{eq:delta-cf}.
This is the concept contrast that the pseudo-outcome decomposition isolates
from the covariate arm once $P(X)$ is shared across $D_{\mathrm{exist}}$
and $D_{\mathrm{new}}$.
The benchmark does not re-estimate a propensity: under the common
covariate table that weight is not where the signal sits.
Gating by the same skill factor as in \eqref{eq:delta-region} produces
$\tau_k^{2,\mathrm{g}}$.
Raw $\tau_k^2$ is retained because a flat local predictor on a noise site
can still sit far from a sharp reference predictor.
The gate is the part of $\Delta_k^{\mathrm{cf}}$ that refuses to call that
gap a primary source.

For a score $s$ and two disjoint collections of sites $\mathcal{A}$ and
$\mathcal{B}$, the pairwise concordance is
\begin{equation}
\label{eq:concordance}
C(s;\mathcal{A},\mathcal{B})
=
\frac{1}{|\mathcal{A}|\,|\mathcal{B}|}
\sum_{a\in\mathcal{A}}\sum_{b\in\mathcal{B}}
\mathbf{1}\{s_a>s_b\}.
\end{equation}
Ties count as failures.
The primary set $\mathcal{P}$ is the five structured sites.
The target for the headline score is
$C(\Delta^{\mathrm{cf}};\mathcal{P},\mathcal{C})>0.8$,
where $\mathcal{C}$ is the set of four clean sites.
The same functional is reported for $\mathcal{P}$ against the noise site,
and separately for the linear and nonlinear subsets of $\mathcal{P}$.
Every entry below is a mean over four seeds, with the worst seed in
parentheses, unless a row is explicitly a per-type mean.

Negative controls use the same windows and never see the label rule.
The domain AUC is the cross-fit AUC of a classifier trained to separate
burn-in rows from recent rows using $X$ alone, once with the forest and
once with logistic regression.
The covariate discrepancy is an RBF MMD on subsamples of size $250$, with
bandwidth equal to the median of the squared Euclidean distances.
A joint classifier two-sample test appends $y$ to $X$ before the same
forest AUC is computed.
Marginal label total variation is the absolute difference of positive rates
between the recent window and the burn-in.

\begin{table}[t]
\centering
\caption{Mean score by site type at $\alpha=1$, abrupt activation, and
$q=0.5$. Four seeds. $\Delta^{\mathrm{cf}}$ is \eqref{eq:delta-cf}.
$\Delta(\mathrm{all})$ is the whole-window score.
$\Delta(\mathrm{high})$ is the upper tail of $X_1$.
Domain AUC is the logistic classifier on $X$ only.}
\label{tab:means}
\small
\begin{tabular}{llrrrrrr}
\toprule
Task & Type & $\Delta^{\mathrm{cf}}$ & $\Delta(\mathrm{all})$ & $\Delta(\mathrm{high})$ & $\tau^2$ & Skill & Domain AUC \\
\midrule
Income & Clean & 0.003 & 0.000 & 0.002 & 0.015 & 0.696 & 0.498 \\
Income & Linear & 0.253 & 0.024 & 0.253 & 0.181 & 0.937 & 0.482 \\
Income & Nonlinear & 0.336 & 0.298 & 0.316 & 0.313 & 0.897 & 0.494 \\
Income & Noise & 0.001 & 0.000 & 0.000 & 0.096 & 0.008 & 0.494 \\
Public Coverage & Clean & 0.001 & 0.000 & 0.001 & 0.014 & 0.364 & 0.503 \\
Public Coverage & Linear & 0.200 & 0.052 & 0.200 & 0.132 & 0.943 & 0.493 \\
Public Coverage & Nonlinear & 0.236 & 0.182 & 0.179 & 0.159 & 0.970 & 0.478 \\
Public Coverage & Noise & 0.003 & 0.002 & 0.003 & 0.046 & 0.089 & 0.528 \\
\bottomrule
\end{tabular}
\end{table}

Table~\ref{tab:means} is the raw material of the concordance claim.
On both tasks, and for every linear site, the maximum is exactly the
upper-tail score: $\Delta^{\mathrm{cf}}=\Delta(\mathrm{high})$.
On Income that tail score is $0.253$ against a whole-window score of
$0.024$.
The whole window is mostly unrewritten rows, and a single Brier average
dilutes a real halfspace by about an order of magnitude.
The noise skill is $0.008$ on Income and $0.089$ on Public Coverage, so
the gate shuts even though the ungated contrast $\tau^2$ on Income noise
($0.096$) exceeds the clean contrast ($0.015$).
Logistic domain AUCs sit on $1/2$ for every type.
The largest is $0.528$, on the Public Coverage noise site.
The forest domain AUC, not shown column by column, has type means between
$0.490$ and $0.514$, and equals $0.514$ on that same noise site.
Neither is evidence of a covariate shift.

\subsection{Concordance}
\label{sec:concordance}

\begin{table}[t]
\centering
\caption{Pairwise concordance \eqref{eq:concordance} at $\alpha=1$, abrupt
activation, and $q=0.5$.
Each cell is the mean over four seeds, with the worst seed in parentheses.
$\mathcal{P}$ is the five structured sites and $\mathcal{C}$ is the four
clean sites.}
\label{tab:concordance}
\small
\begin{tabular}{llcccc}
\toprule
Task & Score & $\mathcal{P}>\mathcal{C}$ & $\mathcal{P}>\mathrm{noise}$ & Linear $>\mathcal{C}$ & Nonlinear $>\mathcal{C}$ \\
\midrule
Income & $\Delta^{\mathrm{cf}}$ & 1.000 {\scriptsize(1.000)} & 1.000 {\scriptsize(1.000)} & 1.000 {\scriptsize(1.000)} & 1.000 {\scriptsize(1.000)} \\
Income & $\Delta(\mathrm{all})$ & 1.000 {\scriptsize(1.000)} & 1.000 {\scriptsize(1.000)} & 1.000 {\scriptsize(1.000)} & 1.000 {\scriptsize(1.000)} \\
Income & $\tau^{2,\mathrm{g}}$ & 1.000 {\scriptsize(1.000)} & 1.000 {\scriptsize(1.000)} & 1.000 {\scriptsize(1.000)} & 1.000 {\scriptsize(1.000)} \\
Income & Pareto rank & 1.000 {\scriptsize(1.000)} & 1.000 {\scriptsize(1.000)} & 1.000 {\scriptsize(1.000)} & 1.000 {\scriptsize(1.000)} \\
Income & $\tau^2$ & 1.000 {\scriptsize(1.000)} & 1.000 {\scriptsize(1.000)} & 1.000 {\scriptsize(1.000)} & 1.000 {\scriptsize(1.000)} \\
Income & Pooled excess & 1.000 {\scriptsize(1.000)} & 0.400 {\scriptsize(0.400)} & 1.000 {\scriptsize(1.000)} & 1.000 {\scriptsize(1.000)} \\
Income & Label TV & 0.500 {\scriptsize(0.100)} & 0.050 {\scriptsize(0.000)} & 0.500 {\scriptsize(0.083)} & 0.500 {\scriptsize(0.125)} \\
Income & Joint C2ST & 0.850 {\scriptsize(0.800)} & 0.350 {\scriptsize(0.200)} & 0.750 {\scriptsize(0.667)} & 1.000 {\scriptsize(1.000)} \\
Income & RF domain AUC & 0.487 {\scriptsize(0.400)} & 0.500 {\scriptsize(0.000)} & 0.521 {\scriptsize(0.333)} & 0.438 {\scriptsize(0.250)} \\
Income & RBF MMD & 0.412 {\scriptsize(0.300)} & 0.500 {\scriptsize(0.000)} & 0.375 {\scriptsize(0.167)} & 0.469 {\scriptsize(0.125)} \\
\midrule
Public Coverage & $\Delta^{\mathrm{cf}}$ & 1.000 {\scriptsize(1.000)} & 1.000 {\scriptsize(1.000)} & 1.000 {\scriptsize(1.000)} & 1.000 {\scriptsize(1.000)} \\
Public Coverage & $\Delta(\mathrm{all})$ & 1.000 {\scriptsize(1.000)} & 1.000 {\scriptsize(1.000)} & 1.000 {\scriptsize(1.000)} & 1.000 {\scriptsize(1.000)} \\
Public Coverage & $\tau^{2,\mathrm{g}}$ & 1.000 {\scriptsize(1.000)} & 1.000 {\scriptsize(1.000)} & 1.000 {\scriptsize(1.000)} & 1.000 {\scriptsize(1.000)} \\
Public Coverage & Pareto rank & 1.000 {\scriptsize(1.000)} & 1.000 {\scriptsize(1.000)} & 1.000 {\scriptsize(1.000)} & 1.000 {\scriptsize(1.000)} \\
Public Coverage & $\tau^2$ & 1.000 {\scriptsize(1.000)} & 1.000 {\scriptsize(1.000)} & 1.000 {\scriptsize(1.000)} & 1.000 {\scriptsize(1.000)} \\
Public Coverage & Pooled excess & 1.000 {\scriptsize(1.000)} & 1.000 {\scriptsize(1.000)} & 1.000 {\scriptsize(1.000)} & 1.000 {\scriptsize(1.000)} \\
Public Coverage & Label TV & 0.613 {\scriptsize(0.350)} & 0.000 {\scriptsize(0.000)} & 0.604 {\scriptsize(0.417)} & 0.625 {\scriptsize(0.250)} \\
Public Coverage & Joint C2ST & 0.950 {\scriptsize(0.850)} & 0.450 {\scriptsize(0.200)} & 0.917 {\scriptsize(0.750)} & 1.000 {\scriptsize(1.000)} \\
Public Coverage & RF domain AUC & 0.438 {\scriptsize(0.350)} & 0.200 {\scriptsize(0.000)} & 0.396 {\scriptsize(0.167)} & 0.500 {\scriptsize(0.250)} \\
Public Coverage & RBF MMD & 0.600 {\scriptsize(0.500)} & 0.200 {\scriptsize(0.000)} & 0.562 {\scriptsize(0.417)} & 0.656 {\scriptsize(0.375)} \\
\bottomrule
\end{tabular}
\end{table}

Table~\ref{tab:concordance} is the primary result.
On both tasks and in every seed,
$C(\Delta^{\mathrm{cf}};\mathcal{P},\mathcal{C})=1$,
and the same equality holds for linear sites alone, for nonlinear sites
alone, and for $\mathcal{P}$ against the noise site.
The whole-window score and the skill-gated contrast $\tau^{2,\mathrm{g}}$
match this ranking at full intensity.
Figure~\ref{fig:concordance} plots the clean-site comparison.
The covariate rows do not.
Random-forest domain AUC and RBF MMD have concordances between $0.38$ and
$0.66$, with worst seeds as low as $0.17$.
Label total variation is consistent with a matched positive rate: its
concordance against clean sites is $0.50$ on Income and $0.61$ on Public
Coverage, and its concordance against the noise site is $0.05$ and $0.00$.
The noise site is the unique marginal shift, and a marginal functional
ranks it first.

The joint classifier two-sample test sees any change in $(X,Y)$.
It separates $\mathcal{P}$ from $\mathcal{C}$ at $0.85$ on Income and $0.95$
on Public Coverage, which clears $0.8$ on average, but it ranks the noise
site as the stronger departure: concordance of $\mathcal{P}$ against noise
is $0.35$ and $0.45$.
On Income the linear-only comparison for the joint test falls to $0.75$.
A test of the joint distribution is not a test for a learnable new concept.

Pooled excess makes the same mistake in a sharper way, and only on the
task where the reference concept is easy.
It is the Brier gap between the forest trained on the entire burn-in and
the local cross-fit model, before size matching and before the skill gate.
On Income its concordance against the noise site is $0.400$ in every seed.
That fraction is exact: both nonlinear sites beat the noise site, and none
of the three linear sites do.
The pooled Income forest is a sharp predictor of $y^{\mathrm{base}}$.
Evaluated on coin-flip labels its Brier rises by $0.084$ on average, while
the same forest evaluated on a linear tail that covers $23\%$ of the window
rises by only $0.033$.
Size matching does not by itself fix the comparison.
On the region that attains $\Delta^{\mathrm{cf}}$, the size-matched excess
on Income noise remains $0.071$, against $0.002$ on clean sites.
Multiplying by the skill factor is what sends the noise score from $0.071$
to $0.001$.
Public Coverage does not reproduce the $0.40$ figure.
There the baseline concept is harder, the pooled excess on noise is
$0.035$, and every structured site still outranks it.
The failure mode is real on the task where the reference fit is strong,
which is the task on which an ungated risk gap is most tempting.

\begin{figure}[t]
\centering
\includegraphics[width=\textwidth]{concordance_cd_vs_clean.png}
\caption{Pairwise concordance against clean sites.
The dashed line is the $0.8$ target.
$\Delta^{\mathrm{cf}}$, the whole-window score, skill-gated $\tau^2$, raw
$\tau^2$, and the Pareto rank clear the line on both tasks.
The joint classifier two-sample test also clears it on average; its failure
is the noise contrast in Table~\ref{tab:concordance}.
Label total variation, the random-forest domain AUC, and RBF MMD do not
clear the line.
Source: abrupt activation, $\alpha=1$, $q=0.5$, four seeds.}
\label{fig:concordance}
\end{figure}

The Pareto comparison uses three axes when the time curve below is
available, and two otherwise.
The axes are $\Delta^{\mathrm{cf}}$, $\tau^{2,\mathrm{g}}$, and the time
jump, defined as the maximum of the curve on endpoints at or after
$t^\star$ minus the maximum on endpoints before $t^\star$.
A site's rank is one plus the number of sites that are at least as large
on every axis and strictly larger on at least one.
Concordance is computed on the negative rank, so a smaller rank is a
higher score.
Mean ranks on Income are $1.875$ for nonlinear sites, $3.250$ for linear
sites, $6.625$ for clean sites, and $6.750$ for the noise site.
On Public Coverage they are $1.375$, $2.667$, $7.188$, and $6.500$.
The undominated front is mostly the nonlinear sites: half of them on
Income and five eighths on Public Coverage.
No clean site and no noise site is undominated.
Linear sites are often dominated by a nonlinear site, because the
nonlinear rewrite covers every row and the linear rewrite covers a tail,
but they remain strictly ahead of every clean site.
That is the sense in which the Pareto order puts primary sources first.
It does not claim that every primary source lies on the front.

\subsection{Time localization}
\label{sec:time}

The curve at site $k$ reuses the burn-in reference and recomputes
\eqref{eq:delta-cf} on a sliding window of four batches, using only the
whole window and the upper tail.
The horizontal coordinate in Figure~\ref{fig:time} is the last batch
included in the window.
Drift begins at batch $8$.
A window of length four is entirely inside the drift only when its last
batch is at least $11$.
The figure follows that arithmetic.
Means for clean sites and for the noise site stay below $0.01$ on both
sides of $t^\star$.
Linear and nonlinear means leave zero between batches $9$ and $10$ and
level off at batch $11$.
On Income the nonlinear plateau is near $0.32$ and the linear plateau is
near $0.22$.
On Public Coverage the two plateaus sit together near $0.18$, which is
consistent with the similar headline scores in Table~\ref{tab:means}.

\begin{figure}[t]
\centering
\includegraphics[width=\textwidth]{time_profile_delta_cf.png}
\caption{Mean $\Delta^{\mathrm{cf}}$ against the last batch in a window of
length four.
The vertical line is $t^\star=8$.
The curve levels off at batch $11$, the first endpoint whose window lies
entirely after the change.
Clean and noise traces remain near zero.
Source: abrupt activation, $\alpha=1$, $q=0.5$, four seeds.}
\label{fig:time}
\end{figure}

Two detectors are applied to each curve without being told $t^\star$.
The online rule calibrates on the first two evaluated endpoints, sets the
threshold at that maximum plus $0.01$, and reports the first later endpoint
that exceeds the threshold.
A hit is an estimate within two batches of $t^\star$.
The derivative rule reports the endpoint of the largest first difference.
Table~\ref{tab:time} shows that the derivative is estimating the onset of
the plateau rather than the onset of the drift.
Its mean absolute error is about three batches on every structured type,
which is one window length.
The online rule is the onset estimator.
It hits every nonlinear site on both tasks and every linear site on Public
Coverage.
Income linear sites, whose tail is the thin one, hit in $0.833$ of
replicates, with mean absolute error $1.92$ batches.
The cost is a false alarm rate of $0.375$ on Income clean sites and
$0.125$ on Public Coverage clean sites: any crossing counts, whether or
not it falls near $t^\star$.
Among the Income clean curves that do cross, the mean absolute error is
$4.17$ batches, so the typical alarm is not a call at the change.
A threshold of $0.01$ above a two-point calibration is enough to rank
sites when the window is the full recent block, and it is not enough to
call a changepoint on a short window without a second screen.
The noise site crosses the threshold in two of four Income replicates and
in one of four Public Coverage replicates.
One Income crossing lands on batch $10$, inside the two-batch window; the
other is six batches later.
The Public Coverage crossing is three batches after $t^\star$ and is not
counted as a hit.

\begin{table}[t]
\centering
\caption{Time localization over four seeds.
Online hit rate is the fraction of sites whose first threshold crossing
lies within two batches of $t^\star$.
False-alarm rate is computed only for clean sites.
Derivative MAE is the absolute error of the largest first difference.}
\label{tab:time}
\begin{tabular}{llcccc}
\toprule
Task & Type & Online hit & False alarm & Online MAE & Derivative MAE \\
\midrule
Income & Clean & 0.062 & 0.375 & 4.17 & 4.06 \\
Income & Linear & 0.833 & --- & 1.92 & 2.83 \\
Income & Nonlinear & 1.000 & --- & 1.62 & 3.00 \\
Income & Noise & 0.250 & --- & 4.00 & 3.25 \\
Public Coverage & Clean & 0.062 & 0.125 & 3.00 & 3.88 \\
Public Coverage & Linear & 1.000 & --- & 1.42 & 3.00 \\
Public Coverage & Nonlinear & 1.000 & --- & 1.50 & 3.00 \\
Public Coverage & Noise & 0.000 & --- & 3.00 & 4.00 \\
\bottomrule
\end{tabular}
\end{table}

\subsection{Concept-group type resolution}
\label{sec:type}

Detecting a primary source does not say whether the new concept is the
linear halfspace or one of the two nonlinear rules.
The dictionary is fixed before looking at site labels.
On a window with outcome $y$, three logistic regressions are cross-fit
against an intercept-only Brier score:
$(Z_1,Z_2)$, the scalar $Z_1 Z_2$, and $(\sin Z_1,\, Z_3^2)$.
Group importance is the drop in Brier.
The call is the maximizing group.
One call uses every row in the recent window.
The other uses only the upper tail of $X_1$, the region in which the
linear rule was planted.
Accuracy is computed on the five structured sites only.
Top-five precision asks a different question: of the five sites with the
largest $\Delta^{\mathrm{cf}}$, how many lie in $\mathcal{P}$.

\begin{table}[t]
\centering
\caption{Type resolution at $\alpha=1$, abrupt activation, and $q=0.5$.
Accuracy is the fraction of structured sites whose maximizing concept group
equals the planted mechanism, averaged over four seeds.
Recall is pooled over both tasks and all four seeds, hence over $24$ linear
sites and $8$ sites of each nonlinear mechanism.
Top-five precision of $\Delta^{\mathrm{cf}}$ is $1$ in every seed on both
tasks.}
\label{tab:type}
\begin{tabular}{llcc}
\toprule
Call & Mechanism or task & Accuracy or recall \\
\midrule
All rows & Income, accuracy & 0.700 \\
All rows & Public Coverage, accuracy & 0.700 \\
Upper tail & Income, accuracy & 0.750 \\
Upper tail & Public Coverage, accuracy & 0.900 \\
\midrule
All rows & Linear, recall & 0.500 \\
All rows & Product, recall & 1.000 \\
All rows & Sine-square, recall & 1.000 \\
Upper tail & Linear, recall & 0.917 \\
Upper tail & Product, recall & 0.375 \\
Upper tail & Sine-square, recall & 1.000 \\
\bottomrule
\end{tabular}
\end{table}

Table~\ref{tab:type} separates detection from typing.
Top-five precision is $1$ in every seed: the five largest values of
$\Delta^{\mathrm{cf}}$ are exactly the five structured sites.
Conditional on that detection, the full-window dictionary recovers both
nonlinear mechanisms in every replicate and recovers the linear halfspace
in half of the linear sites.
The miss is structural.
Outside the tail the linear site still carries $y^{\mathrm{base}}$, and on
Income that baseline outcome is already monotone in hours and schooling.
An additive logistic fit can look competitive with the product and
sine-square groups even when the rewritten slice is a halfspace in
$(Z_1,Z_2)$.
Restricting the call to the upper tail raises linear recall from $0.500$
to $0.917$ and raises Public Coverage accuracy from $0.700$ to $0.900$.
The same restriction lowers product recall from $1$ to $0.375$.
The product rule was planted on every row.
Inside the upper tail the factors are truncated, and the additive group
in $(Z_1,Z_2)$ is no longer dominated by the product scalar.
The sine-square group remains separated on either window.
Type resolution is therefore a property of the region on which the group
is scored, not a single accuracy number.

\subsection{Feature support and one-shot aggregation}
\label{sec:features}

Feature scores are computed on the recent window of each structured site
and compared with that site's planted support by AUC.
Four scores are used.
Shift permutation importance is the drop in the squared concept contrast
when a coordinate is permuted inside five quantile bins of the covariate
it is most correlated with.
Both the local forest and the reference forest are re-evaluated on that
same permuted matrix.
The local forest here is fit on the full recent window; the cross-fit in
\eqref{eq:delta-cf} is not reused.
Predictive permutation importance is the drop in held-out AUC for
predicting the recent label.
The difference of that importance between the recent window and the burn-in
is the fourth score's companion; the fourth score itself is the absolute coefficient vector of an
$\ell_1$-penalized logistic regression on the recent window, with
$C=0.5$.

\begin{table}[t]
\centering
\caption{Mean AUC of a feature score against the planted support.
Linear and product supports are $\{X_1,X_2\}$.
The sine-square support is $\{X_1,X_3\}$.
Averages run over the structured sites of four seeds.}
\label{tab:feature}
\begin{tabular}{llccc}
\toprule
Task & Type & Shift permutation & Predictive permutation & $\ell_1$ coefficient \\
\midrule
Income & Linear & 0.906 & 0.927 & 0.786 \\
Income & Nonlinear & 0.789 & 1.000 & 0.930 \\
Public Coverage & Linear & 0.953 & 0.995 & 0.841 \\
Public Coverage & Nonlinear & 0.985 & 1.000 & 0.879 \\
\bottomrule
\end{tabular}
\end{table}

Table~\ref{tab:feature} says that the planted coordinates are recoverable
from a single drifted site.
Predictive permutation AUC is at least $0.927$.
Shift permutation AUC is at least $0.906$ on linear sites and falls to
$0.789$ on Income nonlinear sites, where the product and the sine-square
rule are functions of coordinates that also predict the baseline income
label.
The $\ell_1$ coefficient is weaker than the forest on linear Public
Coverage ($0.841$) and stronger than shift permutation on Income nonlinear
sites ($0.930$).
None of these AUCs is a site-ranking result.
They are the contents of the vector a site would upload.

Aggregation treats each upload as a point on the simplex.
Negative shift-permutation coordinates are truncated at zero and the vector
is renormalized.
The target vector puts mass on a coordinate proportional to the number of
structured sites that plant it.
Uniform averaging, the coordinate-wise median, inverse-variance weights
with variance estimated from the permutation repeats and floored at
$10^{-4}$, and weights proportional to $\Delta^{\mathrm{cf}}$ are compared
by AUC, precision at the number of positive target coordinates, and
Spearman correlation with the target.
Leaving the noise site out of the uniform average, and averaging only the
five structured sites, are oracles used as ceilings rather than as methods.

On Public Coverage every shift aggregator has AUC at least $0.984$, and
the $\Delta^{\mathrm{cf}}$-weighted average and the average that drops the
noise site both have AUC $1$ and precision $1$.
On Income the uniform shift average has AUC $0.869$ and precision $0.667$.
Weighting by $\Delta^{\mathrm{cf}}$ raises those figures to $0.917$ and
$0.750$.
The predictive permutation upload and the $\ell_1$ upload, averaged
uniformly, reach AUC $0.988$ on Income.
That is not evidence against the concept score.
On Income the planted coordinates are hours, schooling, and age, which are
also the coordinates that predict $y^{\mathrm{base}}$ on a clean site.
A clean site's predictive vector already points at the target.
The shift vector is the one that is trying to upload the change rather than
the baseline regression, and it is the vector that improves when clean and
noise sites are down-weighted by $\Delta^{\mathrm{cf}}$.

Leave-one-site-out influence, defined as one minus the Spearman
correlation between the uniform aggregate and the aggregate with one site
removed, does not identify primary sources.
Mean influence on Income is $0.033$ for clean sites, $0.013$ for linear
sites, $0.018$ for the noise site, and $0.021$ for nonlinear sites.
On Public Coverage every type mean lies between $0.005$ and $0.008$.
The uploaded vectors agree with each other closely enough that deleting one
site barely moves the ranking.
Site-level $\Delta^{\mathrm{cf}}$ is the functional that separates primary
sources from clean sites.
Influence on a simplex average is not.

\subsection{Intensity, tail mass, and gradual activation}
\label{sec:sensitivity}

The main grid holds $\alpha=1$, abrupt activation, and $q=0.5$.
The sensitivity grid varies those three levers on Income only, with three
seeds and the same score, and it does not recompute feature uploads or the
joint two-sample test.
Because the $0.5$ and $0.7$ quantiles of hours worked are the same mass
point at $40$ hours, $q=0.7$ selects the same tail as $q=0.5$ and is
omitted.
The quantile $0.8$ moves the threshold to $45$ hours and leaves about
$17\%$ of rows in the linear tail.
The quantile $0.9$ moves it to $50$ hours and leaves about $8\%$.
Nonlinear sites continue to rewrite every activated row.
Only the linear sites become thinner.

\begin{table}[t]
\centering
\caption{Income concordance $C(\Delta^{\mathrm{cf}};\mathcal{P},\mathcal{C})$
under weaker rewrites and thinner linear tails.
Three seeds.
The worst seed is in parentheses.
Quantile $0.5$ is the main tail, about $23\%$ of rows because of the mass
at $40$ hours.}
\label{tab:sensitivity}
\begin{tabular}{lcccccc}
\toprule
& \multicolumn{2}{c}{$q=0.5$} & \multicolumn{2}{c}{$q=0.8$} & \multicolumn{2}{c}{$q=0.9$} \\
\cmidrule(lr){2-3}\cmidrule(lr){4-5}\cmidrule(lr){6-7}
$\alpha$ & Abrupt & Gradual & Abrupt & Gradual & Abrupt & Gradual \\
\midrule
1.0 & 1.000 {\scriptsize(1.000)} & 1.000 {\scriptsize(1.000)} & 1.000 {\scriptsize(1.000)} & 1.000 {\scriptsize(1.000)} & 0.850 {\scriptsize(0.750)} & 0.867 {\scriptsize(0.800)} \\
0.7 & 1.000 {\scriptsize(1.000)} & 1.000 {\scriptsize(1.000)} & 0.967 {\scriptsize(0.950)} & 0.950 {\scriptsize(0.900)} & 0.817 {\scriptsize(0.750)} & 0.683 {\scriptsize(0.500)} \\
0.4 & 0.783 {\scriptsize(0.650)} & 0.750 {\scriptsize(0.650)} & 0.617 {\scriptsize(0.450)} & 0.600 {\scriptsize(0.450)} & 0.550 {\scriptsize(0.450)} & 0.550 {\scriptsize(0.450)} \\
\bottomrule
\end{tabular}
\end{table}

\begin{figure}[t]
\centering
\includegraphics[width=\textwidth]{sensitivity_concordance.png}
\caption{Income concordance of $\Delta^{\mathrm{cf}}$ against clean sites.
Each panel is one linear-tail quantile.
Cell entries are means over three seeds.
Means stay at or above $0.8$ for $\alpha\in\{0.7,1\}$ down to a linear tail
of about $17\%$ of rows.
They fall below $0.8$ when $\alpha=0.4$, and when a tail of about $8\%$ is
combined with $\alpha=0.7$ and a gradual ramp.
Full intensity on that thin tail stays above $0.8$ on average and not in
every seed; the worst seeds are in Table~\ref{tab:sensitivity}.}
\label{fig:sensitivity}
\end{figure}

Table~\ref{tab:sensitivity} and Figure~\ref{fig:sensitivity} locate the
boundary of the $0.8$ target.
At full intensity the concordance remains $1$ down to a linear tail of
about $17\%$ of rows, whether the change is abrupt or spread over four
batches.
At a tail of about $8\%$, full abrupt intensity falls to mean concordance
$0.850$ with a worst seed of $0.750$, so the mean still clears $0.8$ and
the guarantee across seeds does not.
Gradual activation at that thin tail has the same character
($0.867$, worst seed $0.800$).
Lowering the rewrite probability to $0.7$ preserves the target on the main
tail and on the $17\%$ tail.
Combining $\alpha=0.7$ with the $8\%$ tail and a four-batch ramp drops the
concordance to $0.683$, with a worst seed of $0.500$.
At $\alpha=0.4$ every column is below $0.8$.
The main-tail abrupt entry is $0.783$ with a worst seed of $0.650$.
On that same $\alpha=0.4$ slice, the whole-window score without the tail
scan has concordance $0.433$.
The region maximum is what keeps a weak but structured rewrite near the
target.
It does not keep an arbitrarily weak rewrite above the target.

\subsection{What the benchmark supports}
\label{sec:support}

The benchmark supports four claims, each tied to a table.
First, when the rewrite is complete and the linear tail contains about a
fifth of the rows or more, $\Delta^{\mathrm{cf}}$ orders every structured
site ahead of every clean site and ahead of the noise site, on two real
covariate tables, in every seed
(Table~\ref{tab:concordance}).
Second, the functionals that ignore either the existence of a local concept
or the region in which that concept lives do not.
Label total variation and the covariate arms are near chance against clean
sites.
The joint two-sample test and the pooled Brier excess promote the noise
site, completely so for linear sites under the pooled excess on Income.
Third, the time profile of the same score locates an abrupt change to
within about two batches once the scoring window is taken into account,
with a false-alarm rate on clean Income sites that a deployment threshold
would still have to budget for
(Table~\ref{tab:time}).
Fourth, a pre-specified concept-group dictionary can name the nonlinear
mechanisms from the full window and the linear mechanism from the tail, and
it cannot do both from a single region
(Table~\ref{tab:type}).

The benchmark does not support a claim that leave-one-site-out influence
on an averaged variable-importance vector identifies the primary site.
Those vectors are too similar across sites for the deletion of one upload
to move the rank correlation.
It also does not support a claim at very small intensity.
Below a rewrite probability of one half, or inside a linear tail that
covers less than a tenth of the window, pairwise concordance falls through
$0.8$ on at least one seed.
Both limits are visible in the same design that produces the perfect
separation at full intensity, which is the reason for reporting them beside
the headline rather than as a separate simulation.
